% =============================================================================
% Anatomy of a Design Skill: Which Guidance Improves LLM UI Code, and Where It Acts
% Two-part preprint scaffold (Stage-1 ablation + Stage-2 mechanistic interpretability).
%
% OUTCOME-CONTINGENT DRAFT. Every awaited number is a \numslot{ID} (see SLOT_REGISTRY.md);
% every branch the data could take is a \resultbranch{...}{...}. All figures are
% \synthfig previews built from the real p19 machinery on planted (synthetic) data and
% are unmistakably labelled as such. See BUILD.md for compilation.
%
% Standard packages only (no venue .sty). Single-column preprint.
% =============================================================================
\documentclass[10pt]{article}

\usepackage[utf8]{inputenc}
\usepackage[T1]{fontenc}
\usepackage[margin=0.95in]{geometry}
\usepackage{amsmath,amssymb}
\usepackage{graphicx}
\usepackage{booktabs}
\usepackage{array}
\usepackage{longtable}
\usepackage{multirow}
\usepackage{xcolor}
\usepackage{natbib}
\usepackage[colorlinks=true,linkcolor=linknavy,citecolor=linknavy,urlcolor=linknavy]{hyperref}
\usepackage{enumitem}
\usepackage{caption}
\usepackage[protrusion=true,expansion=false]{microtype}

\bibliographystyle{plainnat}

% ------------------------------------------------------------------ colours ---
\definecolor{linknavy}{HTML}{1A3E6E}
\definecolor{slotbg}{HTML}{FFE9A8}     % awaited-number highlight
\definecolor{slottext}{HTML}{7A5A00}
\definecolor{branchbg}{HTML}{E8EEF6}   % result-branch header
\definecolor{branchrule}{HTML}{4A6FA5}
\definecolor{synthbg}{HTML}{FBE6E6}    % synthetic-figure warning

% =============================================================================
% MACRO SYSTEM FOR OUTCOME-CONTINGENT SLOTS
% =============================================================================

% \numslot{ID} : a highlighted placeholder for a number that awaits real data.
%   Renders [NUM:ID]. Every ID is catalogued in SLOT_REGISTRY.md with its source
%   pipeline artefact (notebook/script + column) and the RUNBOOK session that fills it.
\newcommand{\numslot}[1]{%
  \colorbox{slotbg}{\hspace{-0.15em}\textcolor{slottext}{\footnotesize\ttfamily NUM:#1}\hspace{-0.15em}}%
}

% \resultbranch{condition} ... \endresultbranch : a clearly demarcated conditional
%   block -- the prose that SHIPS if that preregistered outcome obtains. Page-breakable
%   (branches are full paragraphs), with a coloured header and a left rule.
\newenvironment{resultbranch}[1]{%
  \par\vspace{0.3em}\noindent
  \colorbox{branchbg}{\small\textcolor{branchrule}{\textbf{\textsf{OUTCOME}}}\ \textsf{#1}}%
  \par\nobreak\vspace{0.1em}%
  \begingroup\setlength{\leftskip}{1.1em}\small%
}{%
  \par\endgroup\vspace{0.35em}%
}

% \synthfig[width]{Fid}{path}{RUNBOOK-session}{caption} : include a synthetic-preview
%   figure with the MANDATORY provenance prefix. Default width 0.80\linewidth.
\newcommand{\synthfig}[5][0.80\linewidth]{%
  \begin{figure}[tbp]\centering
    \fcolorbox{synthbg}{synthbg}{\parbox{0.985\linewidth}{\centering
      \includegraphics[width=#1]{figures/#3}}}
    \caption{\textbf{[#2 -- SYNTHETIC PREVIEW.]} \textit{Synthetic pipeline-validation
      preview; replaced by real data (RUNBOOK session~#4).} #5}
    \label{fig:#2}
  \end{figure}%
}

% shorthands for notation shared with THEORY.md / PLAN.md (identical symbols)
\newcommand{\vhat}{\hat{v}}
\newcommand{\vlstar}{\hat{v}_{\ell^{\!*}}}
\newcommand{\lstar}{\ell^{*}}
\newcommand{\rhostar}{\rho^{*}}
\newcommand{\POC}{\mathrm{POC}}
\newcommand{\ac}{\mathrm{AC1}}
\newcommand{\alphaK}{\alpha_{K}}
\newcommand{\doop}{\operatorname{do}}
\newcommand{\FULL}{\textsc{full}}
\newcommand{\NEUTRAL}{\textsc{neutral}}
\newcommand{\BEAUTY}{\textsc{beauty1}}
\newcommand{\NOSYS}{\textsc{nosys}}
\newcommand{\LOO}{\textsc{loo}}
\newcommand{\AOI}{\textsc{aoi}}

% compact table helpers
\newcolumntype{C}[1]{>{\centering\arraybackslash}p{#1}}

% =============================================================================
\title{\textbf{Anatomy of a Design Skill:\\[2pt]
Which Guidance Improves an LLM's UI Code, and Where It Acts}}

\author{%
  \normalsize (author list withheld for review)\\
  \normalsize Project~19 --- Two-stage study on \texttt{Qwen2.5-Coder-7B-Instruct}
}
\date{Preprint. Outcome-contingent draft; numeric slots pending the GPU execution phase.}

% =============================================================================
\begin{document}
\maketitle

\begin{abstract}
\noindent
Large language models produce visually generic front-end code by default---the
``purple-slop'' aesthetic---and a short natural-language \emph{design skill} placed in the
system prompt is widely reported to fix it. The effect is folklore: everyone reports it,
nobody has measured it. We measure it, on one open-weights model
(\texttt{Qwen2.5-Coder-7B-Instruct}), in two committed stages that describe the \emph{same}
system. \textbf{Stage~1} is a black-box factorial ablation of a five-component design skill
(color, layout, typography, component patterns, negative constraints), with a
length-matched neutral control that isolates design \emph{content} from prompt mass.
\textbf{Stage~2} is white-box mechanistic interpretability: we locate the skill's signal in
the residual stream, extract a diff-in-means ``clean-design'' direction, and causally verify
it by steering the model with \emph{no design prompt} on held-out prompts, including task
types never seen during development. Every quality claim must move \emph{two independent
signals}: deterministic objective metrics on the rendered DOM and screenshot, and a
style-controlled Bradley--Terry preference from a checklist-anchored VLM judge, gated on
chance-corrected human agreement and statistical power; an effect that moves neither is
reported as null.

We find that the intact skill improves quality over the length-matched neutral prompt by
\numslot{HSKILL-POC} (standardized $\delta$ on the primary objective composite) and is
preferred \numslot{HSKILL-WIN} of the time; the skill's causal weight concentrates in
\numslot{ABS-NEC-TOP}, and no single component alone reproduces the full effect
(sufficiency is distributed). At the representation level, the design signal is
\numslot{ABS-LOC} and steering the extracted direction with no design prompt reproduces
\numslot{STEER-REPRO} of the skill's held-out quality gain.

\begin{resultbranch}{Stage~2 -- a causally sufficient design direction exists}
A single mid-network direction, added with no design prompt, raises held-out UI quality on
both oracle signals, is not reproduced by a norm-matched random control, and---when projected
out while the skill is present---attenuates the skill's own gain; to our knowledge this is the
first causal, steerable direction for the \emph{visual quality} of generated front-end code.
\end{resultbranch}

\begin{resultbranch}{Stage~2 -- the honest null}
No single linear direction steers held-out UI quality under the coherence guardrail: with
tight confidence intervals that rule out under-powering, we conclude that---in this
model---``clean design'' behavior is distributed or not linearly steerable, a precise negative
result that bounds the linear-representation account of aesthetic guidance.
\end{resultbranch}

\noindent\textit{The honest-null option is a first-class outcome, not a fallback: it is
written up with the same rigor as a positive finding.}
\end{abstract}

\section{Introduction}
\label{sec:intro}

Ask a code model for a landing page and you tend to get the same page: a purple-to-indigo
gradient hero, an \texttt{Inter} headline over one big statistic, three centered feature
cards with rounded corners. Practitioners call it ``AI slop'' or ``the purple gradient
website,'' and attribute it to models sampling the high-probability center of their web
training data---the median of a decade of framework tutorials. The widely reported fix is a
short natural-language \emph{design skill}: a few hundred tokens of guidance, dropped into the
system prompt, that tells the model to make deliberate, opinionated choices and to avoid the
slop tells. Anthropic ships such a skill, and its plugin reports hundreds of thousands of
installs \citep{anthropic2025frontendskill,anthropic2025frontendblog}; the community credits
it widely. Yet the effect has never been measured: which parts of the guidance carry the
causal weight, how large the effect is under a defensible quality oracle, and---more
fundamentally---\emph{where inside the model} a few sentences of prose act to change what a
page looks like.

This paper answers two questions about one system, and insists they describe the same system.
\textbf{(Q1) Which guidance matters?} We decompose a design skill into five components and run
a factorial ablation: leave-one-out for necessity, add-one-in for sufficiency, and a
resolution-V fraction for interactions, against a control that holds prompt length constant so
the effect cannot be an artifact of prompt mass. \textbf{(Q2) Where does it act?} We locate the
skill's signature in the residual stream, extract a diff-in-means direction, and test whether
that direction, injected with \emph{no design prompt}, causally reproduces the skill's quality
gain on held-out prompts. Because Stage~2 requires white-box activation access, the subject
model is open-weights and self-hosted, and Stage~1 runs on the \emph{same} model
(\texttt{Qwen2.5-Coder-7B-Instruct}); otherwise the link between ``which instruction matters''
and ``where it lives'' would be severed.

The methodological spine is a two-signal \emph{oracle} we hold to throughout: no quality claim
counts unless it moves (i) deterministic objective metrics computed on the rendered DOM and
screenshot---never raw code strings---and/or (ii) a style-controlled pairwise preference from a
checklist-anchored VLM judge, admitted only after passing a chance-corrected human-agreement
gate and a power gate. This is a direct response to two facts from the literature: objective
interface-aesthetics metrics explain at most about half the variance in human appeal
\citep{miniukovich2015computation,reinecke2013predicting}, so the objective half is necessary
but insufficient; and (V)LM judges are reliable on UIs when checklist-anchored
\citep{zhang2025artifactsbench,jeon2025gfocus} but biased by length and verbosity
\citep{zheng2023mtbench}, so the preference half must be style-controlled and human-validated.
An effect that moves neither signal is reported as null, and a null is a result.

\paragraph{Contributions.} Stated precisely, and credited against the nearest prior work in
Section~\ref{sec:related}:
\begin{enumerate}[leftmargin=1.4em,itemsep=2pt,topsep=2pt]
\item \textbf{The first causal component-ablation of a ``design skill.''} Existing UI
benchmarks score \emph{models or tools}
\citep{si2024design2code,zhang2025artifactsbench,xiao2025designbench}; to our knowledge none
decomposes the \emph{guidance} and attributes quality to its components with a factorial design
and a two-signal oracle. We report which components are necessary, which are sufficient, how
they interact, and how a 391-token skill compares to a five-word ``make it beautiful'' nudge.
\item \textbf{The first activation-level interpretability of aesthetic UI-code generation.}
Activation steering has been applied to refusal \citep{arditi2024refusal}, persona and style
\citep{chen2025persona,konen2024style}, and broad capabilities
\citep{vanderweij2024broad}---but, to our knowledge, no prior work extracts or steers a
direction for the \emph{visual/aesthetic quality} of generated front-end code. We test the
hard, causal bar: does the direction reproduce the gain with no design prompt, on unseen task
types, under a coherence guardrail, with a specificity control and a necessity flip test.
\item \textbf{A prompt-level$\leftrightarrow$representation-level correspondence bridge.} We ask
whether the components the ablation finds most causal map onto separable steering
sub-directions inside the model---do color/layout/typography ablation signatures correspond to
near-orthogonal activation directions that each steer their own metric family? This links a
black-box explanation to a white-box one, a bridge predicted plausible by the geometry of
linear representations \citep{park2024geometry}.
\item \textbf{The honest null as a first-class outcome.} Every research question carries a
preregistered null branch. If no clean steerable direction exists, we report that design
behavior is distributed or not linearly steerable in this model, with confidence intervals
tight enough to distinguish a real null from under-powering---a legitimate, publishable result,
never forced and never buried.
\end{enumerate}

\noindent
We do not claim to invent diff-in-means extraction, style/persona steering, the two-signal UI
oracle, or the objective aesthetics metrics; each is credited in Section~\ref{sec:related}. The
contribution is the synthesis pointed at a new target---design skill $\rightarrow$ UI code
$\rightarrow$ activations---with the causal rigor (both steering arms, a specificity control, an
honest failure surface) the target demands.

\paragraph{Reading map.} Section~\ref{sec:related} positions the work. Section~\ref{sec:design}
gives the study design (research questions, the skill and its decomposition, the corpus, the
confound controls, preregistration). Section~\ref{sec:oracle} defines the two-signal oracle.
Section~\ref{sec:stage1} reports Stage~1 (RQ1--RQ4) with a completed interpretation for every
preregistered outcome. Section~\ref{sec:stage2method} gives the Stage~2 method and
Section~\ref{sec:stage2results}---the pivotal section---its outcome-contingent results.
Sections~\ref{sec:discussion}--\ref{sec:repro} discuss, bound, and reproduce.
Numeric results awaiting the GPU execution phase are shown as \numslot{EXAMPLE} placeholders,
each catalogued to the pipeline artifact and RUNBOOK session that fills it.

\section{Related work}
\label{sec:related}

We situate the study in four literatures and name, in each, the nearest neighbor and the
precise gap we fill.

\subsection{UI-code generation and its evaluation}
A fast-growing line generates front-end code and asks how good it is. \citet{si2024design2code}
(Design2Code) frame screenshot-to-HTML and, crucially for us, evaluate on the \emph{rendered}
page with block-level visual metrics rather than raw code---the machinery we adopt in a
reference-free form. \citet{wu2024uicoder} improve UI code via automated feedback with a hard
render-success/compilability gate; datasets such as WebSight \citep{laurencon2024websight} and
Web2Code \citep{yun2024web2code} scale the paired data. On evaluation, ArtifactsBench
\citep{zhang2025artifactsbench} reports that a checklist-anchored MLLM judge reaches
$\sim$94\% agreement with human/WebDev-Arena rankings; DesignBench \citep{xiao2025designbench},
UI-Bench \citep{uibench2025}, Vision2Web \citep{he2026vision2web}, and WebDev Arena
\citep{lmarena2025webdev} score \emph{models or tools}; G-FOCUS \citep{jeon2025gfocus}
studies robust pairwise UI assessment with position-bias control. Accessibility is measurable
deterministically via axe-core \citep{deque2024axecore} against WCAG~2.1 \citep{w3c2018wcag21}.
\textbf{Gap.} All of these score \emph{a system's output quality}. None decomposes the
\emph{guidance} that produces the output and attributes quality to its components. Our Stage~1
turns the evaluation apparatus of this literature into the response variable of a controlled
experiment on the instruction.

\subsection{Computational aesthetics of interfaces}
Quantifying visual quality has a long history: \citet{ngo2003modelling} define balance,
equilibrium, symmetry, and regularity from element geometry; \citet{miniukovich2015computation}
extend these to real interfaces and---centrally for our oracle---find they explain
$\lesssim$49\% of the variance in perceived aesthetics. \citet{reinecke2013predicting} show
colorfulness (a moderate optimum) and visual complexity (an inverted-U) predict first
impressions with $\sim$50\% ceiling; \citet{hasler2003colorfulness} give the colorfulness
statistic we use. Learned scorers such as UIClip \citep{wu2024uiclip} assess UI design
data-drivenly. \textbf{Gap and consequence.} Because these metrics cap at roughly half the
variance, we treat every objective metric as a \emph{diagnostic}, gate learned scorers behind a
human-correlation check, and make aesthetic \emph{quality} claims ride on the paired preference
signal---the two-signal oracle of Section~\ref{sec:oracle}. We also turn the skill's own
negative constraints into a validated ``purple-slop index,'' making the folk notion of AI slop
a computable quantity.

\subsection{Prompt- and instruction-component analysis}
The methodological shape of Stage~1---compare a behavior with vs.\ without a piece of the
instruction---appears in contrastive-prompt steering \citep{rimsky2024caa,chen2025persona}, and
the confound it must defeat (adding content also adds tokens) is exactly what Chatbot-Arena
style control removes by regressing response length and format out of a Bradley--Terry model
\citep{chiang2024stylecontrol}. Persona-vector extraction pointedly uses \emph{length- and
format-matched} contrastive system prompts \citep{chen2025persona}, the analytic twin of our
length-matched neutral control. We formalize the ablation as a fractional-factorial experiment
\citep{montgomery2017doe}, separating necessity (leave-one-out), sufficiency (add-one-in), and
interactions (a resolution-V fraction). \textbf{Gap.} We are not aware of prior work applying a
factorial design plus a two-signal oracle to the components of a \emph{design} skill for code
generation; the synthesis of DOE, confound-matched controls, and a rendered-output oracle is
new to this target.

\subsection{Activation steering and mechanistic interpretability}
Stage~2 builds on the diff-in-means lineage. \citet{arditi2024refusal} show refusal is mediated
by a single residual-stream direction, extracted by difference-in-means and verified by two
causal arms---activation \emph{addition} and directional \emph{ablation} (weight
orthogonalization); this is our exact skeleton, with ``skill vs.\ length-matched neutral''
replacing ``harmful vs.\ harmless.'' Contrastive Activation Addition \citep{rimsky2024caa} and
ActAdd \citep{turner2023actadd} establish add-a-vector steering; style vectors
\citep{konen2024style} and persona vectors \citep{chen2025persona}---our nearest
neighbors---steer \emph{style/trait} using the mean-over-response-tokens statistic we adopt for
free-form generation, and Persona Vectors do so on the same model family
(Qwen2.5-7B-Instruct), de-risking feasibility. Representation Engineering \citep{zou2023repe}
supplies the PCA/LAT cross-check. The reliability caveats are decisive for how we report:
steering vectors are heterogeneous, with up to $\sim$50\% \emph{anti-steerable} inputs for some
concepts \citep{tan2024analysing} and in/out-of-distribution fragility
\citep{dasilva2025steeringoff}, so we report per-prompt distributions and an anti-steerable
fraction, pre-screen by class separation, and hold out unseen task \emph{types}. On why
diff-in-means rather than sparse autoencoders: AxBench \citep{wu2025axbench}, the largest
head-to-head benchmark, finds diff-in-means beats SAE steering and leads concept detection, and
there is no public SAE suite for Qwen2.5-Coder \citep{deng2026qwenscope}; steering-vector SAE
decompositions can be unfaithful \citep{mayne2024decompose}. We therefore adopt diff-in-means as
primary and demote SAEs to a labeled stretch. We locate with linear probes plus control-task
selectivity \citep{hewitt2019control,alain2017probes} and denoising activation patching
\citep{zhang2024patching,heimersheim2024patching}; we interpret through the
linear-representation hypothesis and superposition
\citep{park2023lrh,park2024geometry,elhage2022superposition}, while keeping non-identifiability
on the table---a ``single direction'' may be one face of a concept cone
\citep{wollschlager2025cones}, and not all features are linear \citep{engels2024notlinear}.
Steering strength is non-monotone \citep{taimeskhanov2026strength}, motivating a swept
coefficient under a KL guardrail; broad-skill \citep{vanderweij2024broad} and conditional
\citep{lee2024cast} steering supply our fallbacks. \textbf{Gap.} An exhaustive sweep of this
literature finds no work extracting or steering a direction for the \emph{visual/aesthetic
quality} of generated front-end code. That is the differentiated Stage~2 contribution, and we
hold it to the causal bar the reliability literature demands.

\section{Study design: problem framing and data}
\label{sec:design}

\subsection{The subject model, locked}
Both stages run on \texttt{Qwen2.5-Coder-7B-Instruct} \citep{hui2024qwencoder} (28 layers, width
$d=3584$). The lock is
load-bearing: Stage~2 needs white-box activation access, so the model is open-weights and
self-hosted, and Stage~1 uses the \emph{same} model so that ``which instruction matters'' and
``where it lives'' describe one system. API models appear only as evaluation judges, never as
the subject. Sampling is fixed across every cell and both engines: temperature $0.7$,
top-$p$ $0.9$, top-$k$ $40$, repetition penalty $1.05$, $\le 4096$ new tokens, seeds
$\{0,\dots,4\}$ (interaction cells $\{0,1,2\}$). Temperature $0.7$ is deliberate: greedy
decoding would collapse seeds to a single output and defeat the distributional analysis.

\subsection{Research questions and decision rules}
Table~\ref{tab:rq} states each research question as \emph{hypothesis $\cdot$ endpoints $\cdot$
decision rule $\cdot$ outcomes}, compressed from the preregistration. The two primary endpoints
recur throughout: the \textbf{Primary Objective Composite} (POC; Section~\ref{sec:oracle}) as
the objective signal, and a \textbf{style-controlled Bradley--Terry utility} as the preference
signal. ``Reliability-gated'' means the judge passed the human-agreement gate; ``powered''
means the contrast met its McNemar pair count for the observed effect. Every RQ has an explicit
honest-null branch (Section~\ref{sec:stage1results}, \ref{sec:stage2results}).

\begin{table}[tbp]
\centering\footnotesize
\caption{Research questions, endpoints, and preregistered decision rules (Table~T5). Holm
controls the family-wise error over the 12-test confirmatory family at $0.05$; BH controls FDR
at $0.10$ for exploratory scans. The full frozen rules are in the preregistration
(Appendix~\ref{app:prereg}).}
\label{tab:rq}
\begin{tabular}{@{}p{0.055\linewidth}p{0.30\linewidth}p{0.35\linewidth}p{0.20\linewidth}@{}}
\toprule
\textbf{RQ} & \textbf{Question / hypothesis} & \textbf{Decision rule} & \textbf{Outcomes} \\
\midrule
-- & H-skill: \FULL$\succ$\NEUTRAL{} on both signals (skill is content, not mass) &
non-null per oracle; report $\delta$ and \%-gap over \NOSYS & works / null \\
RQ1 & Component \emph{necessity}: $\mu(\FULL)-\mu(\LOO\text{-}C_i)>0$ &
$C_i$ necessary if $\delta>0$ and Holm-sig.\ on POC \textbf{or} gated BT CI${>}0$; rank by $\delta$
& necessary (ranked) / null / \emph{negative} \\
RQ2 & Component \emph{sufficiency}: $\mu(\AOI\text{-}C_i)-\mu(\NEUTRAL)>0$, ${<}\FULL$ &
sufficient-partial if Holm-sig./gated; \emph{distributed} if none reaches $30\%$ of the gap &
one-carries / distributed / null \\
RQ3 & \emph{Interactions}: some two-factor interaction ${\neq}0$ &
pair interacts if BH-sig.\ on POC and concordant on ${\ge}1$ other family & additive /
interacting (named) \\
RQ4 & Nudge: $\mu(\FULL)>\mu(\BEAUTY)>\mu(\NEUTRAL)$ &
content if \FULL${-}$\BEAUTY${>}0$ gated; nudge-suffices if \BEAUTY${\ge}70\%$ gap and
\FULL${-}$\BEAUTY{} null & content / nudge / neither \\
RQ5 & Locate: signal peaks in a contiguous mid-band &
band $=$ layers with AUC${\ge}0.80$, selectivity${\ge}0.15$, patch${\ge}25\%$ & localized /
diffuse \\
RQ6 & \textbf{The bar}: steer with \emph{no} design prompt raises held-out quality &
four-conjunct success criterion (Section~\ref{sec:stage2method}) & sufficient (\&/or
necessary) / null \\
RQ7 & Failure surface: generalizes OOD with bounded anti-steerable fraction &
unseen-type $\delta{>}0$ gated and anti-steerable ${<}50\%$ & robust / partial / brittle \\
RQ8 & Correspondence: component sub-vectors near-orthogonal, each steers its family &
mean $|\!\cos|{<}0.3$ and ${\ge}3/5$ signatures match & found / partial / absent \\
\bottomrule
\end{tabular}
\end{table}

\subsection{The independent variable: a five-component design skill}
The skill (frozen text in Appendix~\ref{app:skill}) is decomposed into five components, adapted
from the Anthropic \texttt{frontend-design} skill and community ``purple-slop'' discourse
\citep{anthropic2025frontendskill}: \textbf{C1} color/palette, \textbf{C2}
layout/spacing/hierarchy, \textbf{C3} typography, \textbf{C4} component patterns, \textbf{C5}
negative constraints. A deliberate design choice makes C1--C4 \emph{purely prescriptive} and
consolidates \emph{all} negative constraints in C5, so the components are content-orthogonal:
ablating C5 removes all ``avoid'' guidance, and no ``no purple gradients'' hides inside the
color component. The five components are token-balanced to within $\pm15\%$ of the $\sim$85-token
mean (a build-time assertion using the Qwen tokenizer); the full skill is 391 exact tokens.

\subsection{The 24-cell factorial and the resolution-V fraction}
\label{sec:cells}
Encode each cell by a sign vector $x\in\{\pm1\}^5$ ($+1$ real component, $-1$ its matched
filler; factor map $A\!\dots\!E=C_1\!\dots\!C_5$). We take the $2^{5-1}$ principal fraction with
generator $E{=}ABCD$, defining relation $I{=}ABCDE$, which keeps the 16 corners with an even
number of $-1$'s. This fraction is \emph{resolution~V}: all five main effects and all ten
two-factor interactions are estimable, each aliased only with (assumed-negligible) three- and
four-factor interactions (Appendix~\ref{app:theory}). Six of its 16 runs coincide with named
cells---the all-$+$ run is \FULL, the five single-$+$ runs are the \AOI{} cells---so it
contributes 10 new three-component cells. With four controls, five leave-one-out (\LOO) cells,
and five add-one-in (\AOI) cells, the deduplicated design has $4+5+5+10=\mathbf{24}$ distinct
cells (Table~\ref{tab:cells}). \LOO{} and \NEUTRAL{} lie in the \emph{complementary} fraction
(odd $\#$minus) and enter as replicated paired contrasts, not as factorial-model rows.

The reason necessity (RQ1) and sufficiency (RQ2) are asked separately is exact. Writing the
multilinear factorial model and evaluating at the relevant corners (Appendix~\ref{app:theory},
Eq.~\ref{eq:necsuff}) gives, for component $C_i$,
\begin{equation}
\underbrace{\mathbb{E}[y_{\FULL}]-\mathbb{E}[y_{\LOO\text{-}C_i}]}_{\text{necessity } \mathrm{Nec}_i}
= 2\beta_i + 2\!\!\sum_{j\neq i}\!\beta_{ij} + R,
\qquad
\underbrace{\mathbb{E}[y_{\AOI\text{-}C_i}]-\mathbb{E}[y_{\NEUTRAL}]}_{\text{sufficiency } \mathrm{Suff}_i}
= 2\beta_i - 2\!\!\sum_{j\neq i}\!\beta_{ij} + R,
\label{eq:necsuff}
\end{equation}
so necessity and sufficiency coincide only when a component's interactions vanish; their average
recovers the main effect $2\beta_i$ and their half-difference isolates its total interaction
load. Positive synergies ($\sum_j\beta_{ij}>0$) make a component look \emph{more necessary than
sufficient}---components carry the skill together. This single identity ties RQ1, RQ2, and RQ3
and predicts the directional prior that C5 (negatives) helps more once a palette is present.

\begin{table}[tbp]
\centering\footnotesize
\caption{The 24 distinct cells (Table~T1). $+$ real component, $-$ matched filler,
$\cdot$ not a factorial cell. The 16-run factorial is $\{$\FULL, \AOI-C1..C5, F-123..F-345$\}$;
\LOO{} and \NEUTRAL{} are off-fraction contrast anchors. Headline cells get 5 seeds, interaction
cells 3.}
\label{tab:cells}
\begin{tabular}{@{}llccccccl@{}}
\toprule
\# & cell & C1 & C2 & C3 & C4 & C5 & seeds & tier / role \\
\midrule
1 & \FULL & $+$ & $+$ & $+$ & $+$ & $+$ & 5 & control: intact skill; Stage-2 skill side \\
2 & \NEUTRAL & $-$ & $-$ & $-$ & $-$ & $-$ & 5 & control: length-matched neutral; contrast side \\
3 & \BEAUTY & $\cdot$ & $\cdot$ & $\cdot$ & $\cdot$ & $\cdot$ & 5 & control: 5-word nudge (RQ4) \\
4 & \NOSYS & $\cdot$ & $\cdot$ & $\cdot$ & $\cdot$ & $\cdot$ & 5 & control: bare base; Stage-2 ``no-prompt'' base \\
5--9 & \LOO-C1..C5 & \multicolumn{5}{c}{one $-$, rest $+$} & 5 & necessity (RQ1) \\
10--14 & \AOI-C1..C5 & \multicolumn{5}{c}{one $+$, rest $-$} & 5 & sufficiency (RQ2); in fraction \\
15--24 & F-123..F-345 & \multicolumn{5}{c}{three $+$, two $-$} & 3 & interactions (RQ3); in fraction \\
\bottomrule
\end{tabular}
\end{table}

\subsection{Confound controls as methodological contributions}
Three controls are engineered, not incidental, and we present them as contributions.
\emph{(i) Constant prompt mass.} Every factorial and \LOO{} cell holds 391 exact tokens: a
removed component is replaced \emph{in its original slot} by a token-length-matched filler, so
cells differ only in \emph{content}, never in length or slot order. \NEUTRAL{} is the all-filler
$(-,-,-,-,-)$ corner---a length-matched design-irrelevant prompt, which is why a \FULL$\succ$\NEUTRAL{}
gain is design content and not prompt mass. \emph{(ii) Filler render-inertness.} A filler may
touch only source-level conventions (indentation, comments, naming, declaration order) and is
forbidden by a build-time audit from mandating anything an objective metric could see---semantic
landmarks, \texttt{lang}/\texttt{title}/\texttt{meta}, ARIA, validation, or event handlers---so
fillers cannot themselves move the oracle and bias component effects downward. \emph{(iii)
Constant-mass equalization.} A frozen padding pool deterministically pads or trims each filler
until it matches its component's token count to within $\pm2$ tokens, recorded in the manifest.
A constant output constraint (single-file hermetic HTML) sits in the \emph{user} turn, so it
survives even in \NOSYS, which has no system message.

\subsection{Corpus and the frozen out-of-distribution split}
The corpus is 58 realistic briefs across 10 task types, describing domain and function only
(full table in Appendix~\ref{app:corpus}). A mechanical leakage audit asserts zero whole-word
matches of a frozen list of style adjectives and design-family words over the briefs, so the
system-prompt treatment is the only design signal in context. The split is frozen by id:
\textbf{40 dev} (8 development types $\times$ 5), \textbf{8 held-out seen-type} (one per
development type), and \textbf{10 held-out unseen-type}---settings-panel and admin-data-table,
two control-/information-dense layouts \emph{entirely absent} from dev. The unseen types are the
out-of-distribution generalization test for steering: control-dense layouts unlike the
marketing-oriented dev types stress the skill's guidance in a regime it was not tuned on. A
\emph{generation unit} is one $(\text{cell},\text{prompt},\text{seed})$ triple; Stage~1
comprises 4{,}000 dev units (all 24 cells $\times$ 40) plus 630 held-out units (reduced set),
totaling 4{,}630.

\subsection{Preregistration}
The cell table, the 58-brief split, the POC formula and its two conditionals, the confirmatory
12-test family, the reliability and power gates, the Stage-2 success criterion, and the null
rule are frozen in a preregistration (Appendix~\ref{app:prereg}) that is git-tagged before any
GPU run. A pilot gate (2 prompts $\times$ \{\FULL, \NOSYS\} $\times$ 3 seeds) must show a visible
skill effect before the generation budget is spent; the sole sanctioned model-swap trigger is a
pilot that still shows no effect after prompt-format fixes.

\section{The two-signal oracle}
\label{sec:oracle}

Every quality claim---an ablation delta or a steering delta---must move two independent signals.
This section defines them, the null rule that binds them, and the power that backs them. The
oracle is not a scoring convenience; it is the paper's evidentiary standard, and it is
preregistered.

\subsection{Objective signal: deterministic metrics on the rendered output}
All objective metrics are computed on the \emph{rendered} DOM and screenshot from a pinned,
offline, animation-disabled headless Chromium at $1440\times900$---never on raw code
strings---so they are deterministic and unfoolable. They are reported as a \emph{vector} in six
families: \textbf{V} validity/hermeticity (render success, external-fetch count, overflow and
overlap counts, axe-core violations), \textbf{A} accessibility/contrast, \textbf{L} layout,
\textbf{T} typography, \textbf{C} color, \textbf{X} complexity. Two archetypes fix the flavor.
Per text node we compute the WCAG contrast ratio between foreground luminance $L_1$ and
background $L_2$,
\begin{equation}
\mathrm{CR}=\frac{L_1+0.05}{L_2+0.05}\in[1,21],\qquad
L=0.2126\,R_{\mathrm{lin}}+0.7152\,G_{\mathrm{lin}}+0.0722\,B_{\mathrm{lin}},
\label{eq:wcag}
\end{equation}
(sRGB-linearized channels), reporting the fraction of nodes below $4.5{:}1$. And the skill's own
C5 constraints become a validated \textbf{purple-slop index}
\begin{equation}
\mathrm{PSI}=0.30\,\mathrm{hueFrac}[260^\circ,290^\circ]+0.30\,\mathrm{gradientBg}
+0.25\,\mathrm{InterRobotoShare}+0.15\,\mathrm{centeredHero},
\label{eq:psi}
\end{equation}
each subcomponent in $[0,1]$; the skill tells us exactly what ``slop'' is, so we measure exactly
that. Layout uses Ngo balance/equilibrium/symmetry and an alignment-regularity count; typography
scores modular-scale adherence; color uses Hasler colorfulness $M$ (Appendix~\ref{app:theory}).

\paragraph{The Primary Objective Composite.} To have \emph{one} confirmatory objective number per
generation while honoring ``do not aggregate raw metrics,'' we preregister a mean of oriented,
$z$-scored core metrics,
\begin{equation}
\POC(\mathbf{x})=\tfrac{1}{7}\big[
z(\texttt{align\_reg})+z(\texttt{whitespace})+z(\texttt{type\_scale})+z(\texttt{contrast\_pass})
-z(\texttt{overflow\_overlap})-z(\mathrm{PSI})-z(|M-c^{*}|)\big],
\label{eq:poc}
\end{equation}
$z(\cdot)$ over all dev generation units, oriented so higher is better; the colorfulness term is a
distance to a moderate-optimum target $c^{*}$, not ``more is better.'' Render-failed units are
excluded from POC (a non-rendering page is a validity null, counted in family~V). Two conditionals
are frozen \emph{with} the endpoint, not applied post hoc: (i) if PSI fails its admission gate
(below) the POC drops the $-z(\mathrm{PSI})$ term and every confirmatory decision uses the 6-term
variant; (ii) $c^{*}$ is the median $M$ of human-preferred pages if $\ge30$ exist, else the median
$M$ of dev \FULL-cell generations. The branch that fires is written to the manifest \emph{before}
the first confirmatory test.

\paragraph{Oracle validation gates.} Before any metric is used as evidence: PSI is admitted only
if Spearman $\rho(\mathrm{PSI},\text{human ``AI-generated'' rating})\ge0.4$; on the human subset we
observe $\rho=\numslot{PSI-RHO}$ (Figure~\ref{fig:F10}), so the \numslot{POCVAR} POC is used and
$c^{*}=\numslot{CSTAR}$. The learned UIClip scorer is admitted as \emph{auxiliary} only if
$\rho(\text{UIClip},\text{human win-rate})\ge0.3$ (observed $\numslot{UICLIP-RHO}$); it never enters
the POC or a confirmatory decision, because a learned scorer is not deterministically auditable and
carries a mobile-distribution shift.

\synthfig[0.44\linewidth]{F10}{F10_preview.pdf}{4}{Does PSI track human ``AI-slop'' judgment? The
admission gate ($\rho\ge0.4$) decides whether the confirmatory POC keeps the $-\mathrm{PSI}$ term
(7-term) or drops it (6-term). Axes and the gate annotation are the final layout; the scatter is
planted.}

\subsection{Preference signal: style-controlled Bradley--Terry with a validated judge}
The second signal is a pairwise preference from a checklist-anchored VLM judge (primary:
Gemini~2.5~Flash, a different family from Qwen to avoid self-enhancement; GPT-4o audits 15\%).
Each pair is judged in \emph{both} presentation orders and counts as a winner only if the verdict
is order-consistent, else a tie---a position-bias control. We fit a \textbf{style-controlled
Bradley--Terry} model with a Davidson tie term,
\begin{equation}
P(i\succ j)=\sigma\!\big((\beta_i-\beta_j)+\gamma^{\top}(s_i-s_j)\big),\quad
s=(\log\text{DOM},\ \log\text{tok},\ \#\text{elem},\ \text{colorfulness},\ \text{text-density}),
\label{eq:bt}
\end{equation}
reference $\beta_{\NEUTRAL}=0$, cluster-bootstrap CIs over prompts ($B{=}2000$). The covariates are
judge-bias channels a designer would never score by (node count, token length); omitting them
biases the estimated skill effect by $\gamma^{\top}(\bar s_{\FULL}-\bar s_{\NEUTRAL})$, the
analysis-level twin of the design-level length-matched control. One covariate is flagged as
partially causal: colorfulness \emph{is} a design lever, so regressing it out could absorb a
genuine C1 effect. We therefore report both the controlled and uncontrolled $\hat\beta$, recover
the legitimate color effect through the objective channel (the POC $|M-c^{*}|$ term and the RQ8
color-family steering signature, neither routed through the judge), and if the two $\hat\beta$
disagree in sign we report the contrast as \emph{style-confounded} rather than resolved.

\paragraph{The reliability gate and the kappa paradox.} The preference signal counts only if the
judge tracks a human. Sean labels $N{=}100$ stratified pairs in both orders; we compute
judge--human Krippendorff $\alphaK$ (ordinal) and Gwet $\ac$ with bootstrap CIs. Because the skill
usually wins, the marginal is skewed, and a skewed marginal crushes $\alphaK$ even at high raw
agreement---the kappa paradox. The two coefficients differ only in their chance model,
\begin{equation}
\alphaK=1-\frac{D_o}{D_e},\qquad
\ac=\frac{p_o-p_e^{\gamma}}{1-p_e^{\gamma}},\quad
p_e^{\gamma}=\tfrac{1}{q-1}\textstyle\sum_k \pi_k(1-\pi_k),
\label{eq:agree}
\end{equation}
and $p_e^{\gamma}\to0$ as one category dominates, so $\ac$ stays near $p_o$ where $\alphaK$
collapses (a worked example: at $90\%$ agreement, skew moves $\alphaK$ from $0.80$ to $0.44$ while
$\ac$ holds at $0.88$; Appendix~\ref{app:theory}). The preregistered gate is therefore
\emph{disjunctive}:
\[
\textbf{pass}\iff \alphaK\ge0.667\ \ \textbf{or}\ \ \big(\ac\ge0.80\ \text{under demonstrable skew, }\pi_{\max}\gtrsim0.80\big).
\]
The skew qualifier stops $\ac$ from being a loophole in the balanced case (where the two agree). We
observe $\alphaK=\numslot{REL-ALPHA}$, $\ac=\numslot{REL-AC1}$, $\pi_{\max}=\numslot{REL-PIMAX}$
(Figure~\ref{fig:F11}); the gate is \numslot{REL-GATE}. These thresholds are convention, not
derived optima, and we report the CIs so a reader may apply a stricter bar.

\synthfig[0.44\linewidth]{F11}{F11_preview.pdf}{5}{Can we trust the VLM judge? $\alphaK$, $\ac$, and
the marginal skew $\pi_{\max}$ with bootstrap CIs against the two gate thresholds. The disjunctive
gate admits $\ac$ only under demonstrable skew; the bars are the kappa-paradox exemplar (planted),
the final panel plots the real judge--human numbers.}

\subsection{The null rule (verbatim) and power}
The two signals are bound by a rule we quote exactly as preregistered:
\begin{quote}\itshape
An ablated component's effect, or a steering direction's effect, is reported as null unless it
moves an objective-metric endpoint (POC Holm-significant, or $\ge2$ concordant BH-significant
families) OR a preference delta that passes both the reliability gate and the power gate. Aesthetic
claims not backed by the objective half are inadmissible.
\end{quote}
The disjunction is a deliberate, bounded loosening for the necessity/sufficiency \emph{screen}: with
each channel controlled near $0.05$, the false-non-null rate is $\lesssim0.10$ by a union bound,
which is safe because the dominant risk is a \emph{false null} (missing a real component) and the
objective channel is deterministic. The Stage-2 headline instead uses a conjunction of four checks,
whose joint false-positive rate falls below any single $\alpha$ (Section~\ref{sec:stage2method}).

Preference deltas are paired, so we size them with McNemar's test. Among decisive (order-consistent)
pairs the winner is tested against a fair coin; the required decisive count at alternative split
$\theta_1$ is
\begin{equation}
N_{\mathrm{dec}}\approx\frac{\big(z_{1-\alpha/2}\sqrt{\theta_0(1-\theta_0)}+z_{1-\beta}\sqrt{\theta_1(1-\theta_1)}\big)^2}{(\theta_1-\theta_0)^2},
\quad\theta_0=\tfrac12,
\label{eq:power}
\end{equation}
giving $\approx200$ decisive pairs for a $60/40$ split, $\approx80$ for $65/35$, $\approx780$ for
$55/45$ (at $\alpha{=}0.05$, power $0.80$). The design allocates 200 pairs to \FULL--\NEUTRAL{} and
$\sim$120 to each necessity edge accordingly. Small expected effects (some \AOI{} cells) are lightly
powered on \emph{preference} on purpose---their non-null status can still be established through the
fully powered objective channel, exactly what the two-signal rule permits. The power gate is an
admissibility screen, not an inference: a contrast counts only if it achieved the decisive pairs
Eq.~\ref{eq:power} requires \emph{at the observed} split; an under-powered delta is never
reinterpreted as a null. We observe a tie rate of \numslot{POW-TIE}, \numslot{POW-FN} decisive pairs
on \FULL--\NEUTRAL, \numslot{POW-NEC} per necessity edge, and \numslot{POW-UNDER} under-powered
confirmatory edges (Table~T10, Appendix~\ref{app:extra}).

\section{Stage 1 results: which guidance matters}
\label{sec:stage1}
\label{sec:stage1results}

All Stage-1 objective deltas are contrasts on the POC via a mixed-effects model
$y_{cps}=\mu+\tau_c+u_p+v_{cp}+\epsilon_{cps}$ with fixed cell effect $\tau_c$, random prompt
intercept $u_p$, and random cell-by-prompt slope $v_{cp}$; this is what separates the study from
per-generation $t$-tests, because five seeds on one prompt are not five independent points---with a
modest intraclass correlation the effective $n$ is roughly halved (Appendix~\ref{app:theory}).
Preference deltas are style-controlled Bradley--Terry utilities (Eq.~\ref{eq:bt}). Confirmatory
contrasts are Holm-corrected over the 12-test family, separately on each signal; the per-component
estimates are collected in Table~\ref{tab:s1effects} and visualized in Figure~\ref{fig:F1}.
For each research question we give the completed interpretation the paper will carry under
\emph{every} preregistered outcome; on real data, exactly one branch per question is kept.

\subsection{Does the skill work, and is it content?}
The headline contrast is \FULL{} vs.\ \NEUTRAL. Because \NEUTRAL{} is length-matched, a gain is
design \emph{content}, not prompt mass. We measure $\numslot{HSKILL-POC}$ standardized $\delta$ on
the POC and a \FULL{} win-rate of $\numslot{HSKILL-WIN}$ with BT utility $\numslot{HSKILL-BT}$; the
prompt-mass control \NEUTRAL{} vs.\ \NOSYS{} moves the POC by $\numslot{NEU-NOSYS}$, and \FULL{}
closes $\numslot{HSKILL-NOSYS-PCT}$ of the \NOSYS{} gap.

\begin{resultbranch}{works}
The design skill improves UI quality on both oracle signals: \FULL$\succ$\NEUTRAL{} at
$\numslot{HSKILL-POC}$ POC and $\numslot{HSKILL-WIN}$ preference, both surviving Holm. Because the
gain is measured against a token-matched, design-irrelevant prompt, it is attributable to the
skill's \emph{content} rather than to added prompt mass---the near-zero \NEUTRAL--\NOSYS{} contrast
($\numslot{NEU-NOSYS}$) confirms that filler tokens alone do essentially nothing. The rest of
Stage~1 is well-posed: there is a real effect to decompose.
\end{resultbranch}
\begin{resultbranch}{null}
\FULL$\approx$\NEUTRAL{} on both signals with a tight CI around zero: the model cannot make use of a
structured design skill at all. This outcome would have been caught at the pilot gate and escalated
to the recorded model fallback; were it to survive to the full run, we would report it honestly and
scope the remainder of Stage~1 to ``no measurable skill effect on this model,'' a negative result
about the model's instruction-following rather than about design guidance in general.
\end{resultbranch}

\subsection{RQ1 --- component necessity}
Necessity is the paired contrast \FULL{} vs.\ \LOO-$C_i$: does removing $C_i$ from the intact skill
degrade quality? Per-component objective deltas are in Table~\ref{tab:s1effects} and the twin-panel
forest plot (Figure~\ref{fig:F1}); the directional prior is that C5 (negatives) and C1 (color) punch
above their token weight.

\begin{resultbranch}{necessary (ranked)}
Removing certain components significantly degrades quality; ranked by standardized $\delta$, the
skill's weight concentrates in \numslot{NEC-RANK}, with \numslot{NEC-NSURV} of five components
surviving Holm on the POC. The largest necessity effects are C5 (negatives, $\numslot{NEC-C5}$;
BT utility $\numslot{NEC-C5-BT}$) and C1 (color, $\numslot{NEC-C1}$; BT $\numslot{NEC-C1-BT}$),
confirming the prior that the skill's leverage lies disproportionately in \emph{telling the model
what not to do} and \emph{pinning a palette}---precisely the components whose token weight is
smallest relative to their effect. The ordering replicates on held-out prompts (top component
$\numslot{NEC-REPL-TOP}$), so it is not a dev artifact. Per Eq.~\ref{eq:necsuff}, components with
large necessity but smaller sufficiency (below) carry positive interaction load---they help most in
the presence of the others.
\end{resultbranch}
\begin{resultbranch}{null}
No leave-one-out contrast degrades quality: every \FULL--\LOO-$C_i$ CI brackets zero
(Table~\ref{tab:s1effects}), and the power gate certifies these as tight nulls, not under-powering.
The reading is that no single component is individually necessary---the skill works only as a whole,
with each component substitutable by the others' combined effect. This makes RQ2/RQ3 the load-bearing
questions: a skill that is robust to any single deletion is either highly redundant or purely
interactive.
\end{resultbranch}
\begin{resultbranch}{negative}
Removing a component \emph{improves} quality: at least one \FULL--\LOO-$C_i$ contrast is significantly
\emph{negative} (\LOO{} beats \FULL). This is a genuine, honestly reported finding---an
over-constraining instruction that, in this model, does more harm than good (a plausible candidate is
an over-specified typography or pattern directive that pushes the model into awkward choices). The
practical lesson for skill design is that more guidance is not monotonically better; we report which
component and by how much, and recommend its revision.
\end{resultbranch}

\subsection{RQ2 --- component sufficiency}
Sufficiency is \AOI-$C_i$ vs.\ \NEUTRAL: does $C_i$ \emph{alone} lift quality, and does any single
component reach \FULL? The necessity$\times$sufficiency scatter (Figure~\ref{fig:F2}) reads the two
questions jointly against the $y{=}x$ line where interactions vanish. The largest single-component
lift reaches $\numslot{SUFF-MAXPCT}$ of the \FULL--\NEUTRAL{} gap.

\begin{resultbranch}{distributed}
Every component alone lifts quality above \NEUTRAL{} (\numslot{SUFF-NSURV} of five Holm-significant),
but none reaches the preregistered $30\%$ of the \FULL--\NEUTRAL{} gap---the largest,
\numslot{SUFF-TOP}, attains only $\numslot{SUFF-MAXPCT}$. Sufficiency is therefore \emph{distributed}:
no single instruction carries the skill, and the full effect emerges from the components together.
Read with RQ1 through Eq.~\ref{eq:necsuff}, components that are more necessary than sufficient (above
the $y{=}x$ line in Figure~\ref{fig:F2}) reveal positive synergy---the whole exceeds the sum of parts.
\end{resultbranch}
\begin{resultbranch}{one-carries}
A single component reaches within $30\%$ of \FULL{} on its own: \numslot{SUFF-TOP} attains
$\numslot{SUFF-MAXPCT}$ of the gap, and no other component approaches it. The skill is, to first order,
that one instruction; the remainder is polish. This is a strong, compressible result---most of a
391-token skill's benefit is reproducible with one component---and it sharpens the Stage-2 target: the
extracted direction should correspond most closely to that component's activation signature.
\end{resultbranch}
\begin{resultbranch}{null}
No component alone beats \NEUTRAL{} (all \AOI-$C_i$ CIs bracket zero) even though the intact skill
does (H-skill). Components are effective only in combination: sufficiency is strictly emergent. This
pairs with the RQ3 interaction story---if no main effect is individually sufficient yet the skill
works, the effect lives in the interactions, and the factorial block is where it is found.
\end{resultbranch}

\subsection{RQ3 --- interactions and additivity}
The 16-run factorial estimates all ten two-factor interactions cleanly (each aliased only with an
assumed-negligible three-factor term). The 2FI heatmap is Figure~\ref{fig:F3}; the directional prior
is a positive C1$\times$C5 (negatives help more once a palette is set) and a mild negative C2$\times$C3.

\begin{resultbranch}{interacting}
\numslot{INT-NSURV} two-factor interactions survive BH and are concordant on a second metric family.
The largest is \numslot{INT-MAXPAIR} ($\numslot{INT-MAX}$); as predicted, C1$\times$C5 is
super-additive ($\numslot{INT-C1C5}$)---negative constraints deliver more once a deliberate palette is
specified, because ``avoid purple gradients'' bites hardest when a replacement palette exists---while
C2$\times$C3 is mildly sub-additive ($\numslot{INT-C2C3}$). The components are not independent knobs;
the skill is a small interacting system, and this is the quantitative content of the
necessity$>$sufficiency asymmetry in Eq.~\ref{eq:necsuff}.
\end{resultbranch}
\begin{resultbranch}{additive}
No two-factor interaction survives BH: the component effects simply add, and the skill behaves as a
\emph{checklist} whose value is the sum of its independent items. This is itself informative for skill
design---components can be added, removed, or reordered without surprise---and it means necessity and
sufficiency coincide (Eq.~\ref{eq:necsuff} with $\sum_j\beta_{ij}\approx0$), which Figure~\ref{fig:F2}
would show as points on the $y{=}x$ line.
\end{resultbranch}
\begin{resultbranch}{dominated by a single main effect}
One component's main effect swamps all interactions and most other mains: the skill is effectively one
lever (\numslot{SUFF-TOP}), with the remaining components contributing little independently or
interactively. The paper's Stage-1 story then collapses to ``the skill is essentially component $X$,''
and Stage~2 asks whether $X$'s activation signature is the whole design direction.
\end{resultbranch}

\subsection{RQ4 --- do five words match 391 tokens?}
We compare the intact skill to a five-word nudge (\BEAUTY: ``Make it beautiful and well-designed.'')
and to \NEUTRAL, testing the ordered contrasts. \BEAUTY{} reaches $\numslot{NUDGE-BPCT}$ of the
\FULL--\NEUTRAL{} gap.

\begin{resultbranch}{content}
The ordering \FULL$>$\BEAUTY$>$\NEUTRAL{} holds with \FULL$-$\BEAUTY${>}0$ ($\numslot{NUDGE-FB}$ POC;
\FULL{} preferred over \BEAUTY{} at $\numslot{NUDGE-FBWIN}$): a nudge helps ($\numslot{NUDGE-BN}$ over
\NEUTRAL) but structured guidance is load-bearing and not substitutable by a one-liner. What the
391 tokens buy beyond ``make it beautiful'' is the \emph{specific} content---the palette discipline
and the negative constraints RQ1 finds most necessary---which a generic exhortation cannot supply.
\end{resultbranch}
\begin{resultbranch}{nudge}
\BEAUTY{} reaches $\ge70\%$ of the gap ($\numslot{NUDGE-BPCT}$) and \FULL$-$\BEAUTY{} is null: a
five-word nudge captures most of the skill's effect. This is a striking, cheap result---the model
already ``knows'' how to make a better page and mostly needs \emph{permission} to leave the
high-probability center, not 391 tokens of instruction. It reframes the design skill as an
activation trigger rather than a knowledge transfer, which is exactly the hypothesis Stage~2 tests
mechanistically.
\end{resultbranch}
\begin{resultbranch}{neither}
\BEAUTY$\approx$\NEUTRAL: a vague nudge does nothing, and only the specific structured skill helps (if
H-skill fired). The model needs concrete, itemized guidance or none at all---consistent with a picture
in which ``beautiful'' is too diffuse a target to move the output distribution, whereas named colors,
type scales, and prohibitions are actionable.
\end{resultbranch}

\subsection{RQ7 (black-box slice) --- does the skill effect generalize across task types?}
Before any steering, we report the \FULL--\NEUTRAL{} POC delta \emph{per task type}, including the two
held-out unseen types (settings-panel, admin-data-table). The mean seen-type delta is
$\numslot{BB-SEEN}$ and the mean unseen-type delta is $\numslot{BB-UNSEEN}$. A skill effect that holds
on the information-dense unseen types is evidence the guidance encodes transferable design content, not
a marketing-page idiom; a collapse on those types localizes the skill's benefit to the layouts it was
written for and sets expectations for the Stage-2 out-of-distribution steering test.

\begin{table}[tbp]
\centering\footnotesize
\caption{Stage-1 per-component effects on the POC (standardized $\delta$, 95\% cluster-bootstrap CI),
necessity (\FULL$-$\LOO-$C_i$) and sufficiency (\AOI-$C_i-$\NEUTRAL). Interactions and preference
utilities are reported in Figures~\ref{fig:F1}--\ref{fig:F3}. Every entry awaits the GPU run.}
\label{tab:s1effects}
\begin{tabular}{@{}lccl@{}}
\toprule
component & necessity $\delta$ (\FULL$-$\LOO) & sufficiency $\delta$ (\AOI$-$\NEUTRAL) & prior \\
\midrule
C1 color        & \numslot{NEC-C1} & \numslot{SUFF-C1} & large \\
C2 layout       & \numslot{NEC-C2} & \numslot{SUFF-C2} & medium \\
C3 typography   & \numslot{NEC-C3} & \numslot{SUFF-C3} & medium \\
C4 patterns     & \numslot{NEC-C4} & \numslot{SUFF-C4} & near-null \\
C5 negatives    & \numslot{NEC-C5} & \numslot{SUFF-C5} & largest \\
\bottomrule
\end{tabular}
\end{table}

\synthfig[0.61\linewidth]{F1}{F1_preview.pdf}{5}{Which components carry the skill? Standardized
$\delta$ vs.\ reference for each \LOO{} (necessity), \AOI{} (sufficiency), and \BEAUTY{} cell, with
Holm CIs; filled markers survive Holm. The twin-panel POC/BT layout is final; the effect sizes are
planted (C5, C1 largest; C4 near-null by construction).}

\synthfig[0.39\linewidth]{F2}{F2_preview.pdf}{5}{Necessity $\times$ sufficiency. Points above the
$y{=}x$ line are more necessary than sufficient---positive interaction load (Eq.~\ref{eq:necsuff}).
Planted coordinates; the real figure places the five components by their measured effects.}

\synthfig[0.42\linewidth]{F3}{F3_preview.pdf}{5}{Two-factor interaction coefficients (factorial
MixedLM), BH-significant cells starred. Planted C1$\times$C5 super-additivity and C2$\times$C3
sub-additivity illustrate the encoding; real signs and magnitudes replace them.}

\section{Stage 2 method: locate, extract, steer, verify}
\label{sec:stage2method}

Stage~2 asks where the skill acts and whether that locus is causal. All activation work runs on
Hugging Face \texttt{transformers} with a custom hook module (vLLM cannot inject mid-decode); the
residual stream is the layer-output hidden state, and every comparison is within-engine---we never
compare a vLLM baseline to an HF steered output. We reuse the Stage-1 \FULL{} and \NEUTRAL{} dev
generations (200 matched pairs) as the extraction corpus, capturing the mean residual over response
tokens by teacher-forced forward pass.

\paragraph{Why the residual stream, and why diff-in-means.} Qwen2.5-Coder is a pre-norm decoder:
each block reads a LayerNorm of the stream and writes back additively, so the residual stream is a
linear channel where a concept can be \emph{read} by projection and \emph{written} by addition. The
clean-design direction at layer $\ell$ is the difference of response-mean residuals,
\begin{equation}
v_\ell=\bar h_\ell^{(\FULL)}-\bar h_\ell^{(\NEUTRAL)},\qquad \vhat_\ell=v_\ell/\lVert v_\ell\rVert.
\label{eq:diffmean}
\end{equation}
The contrast against the \emph{length-matched} neutral is what makes $v_\ell$ ``design content''
rather than ``instruction-following mode.'' We steer with this un-whitened vector rather than a probe
direction on purpose: for equal-covariance Gaussians the optimal \emph{discriminant} is
$w^{*}=\Sigma^{-1}(\mu_1-\mu_0)$, ideal for \emph{reading} but pathological for \emph{writing}---it
amplifies low-variance nuisance directions and pushes activations off-distribution---whereas
diff-in-means is the empirical mean shift that lands on the data manifold and cancels nuisance
directions common to both sides (Appendix~\ref{app:theory}). We cross-check $\vhat_\ell$ against the
PCA-of-paired-differences top component and the probe weight, reporting their cosines.

\paragraph{Step 1 --- stability and localization (RQ5).} Extraction adequacy is a running-mean cosine
plateau: as $n$ grows to 200 we require $\Delta\cos<0.01$ over the last $20\%$ (observed
$\numslot{STAB-DCOS}$ at $n=\numslot{STAB-PLATN}$; Figure~\ref{fig:F13}). Localization triangulates
three signals per layer: a linear probe's AUC, its control-task \emph{selectivity}
$\mathrm{AUC}^{\text{real}}-\mathrm{AUC}^{\text{control}}$ (a probe can memorize; selectivity guards
against it), and denoising activation-patching \%-recovered on a design-token logit-diff proxy
(denoising is robust to self-repair). The band is the contiguous layers with AUC$\ge0.80$,
selectivity$\ge0.15$, and patch$\ge25\%$ (Figure~\ref{fig:F5}).

\begin{resultbranch}{RQ5 localized}
The three signals concur on a contiguous mid-band, \numslot{LOC-BAND}: probe AUC peaks at
$\numslot{LOC-AUCPEAK}$ (layer $\numslot{LOC-PEAKLAYER}$) with selectivity $\numslot{LOC-SELPEAK}$ and
patching recovering $\numslot{LOC-PATCHPEAK}$ of the clean--corrupted gap. Decodability, selectivity,
and causal patching agreeing at the same depth is the clean case; these layers become the Stage-2
focus band. Decodability alone would be weak evidence---a probe can read a feature the model does not
use---so the patching agreement is what licenses moving from correlation to a causal candidate.
\end{resultbranch}
\begin{resultbranch}{RQ5 diffuse}
No band meets all three thresholds simultaneously---for instance, probe AUC is high across many layers
but selectivity or patching never clears its bar, indicating the signal is decodable but not localized
to a causal site. We widen to the single peak layer, flag the representation as diffuse, and pre-commit
to multi-layer injection in the steering step. A diffuse locate is not yet a null: it changes the
steering \emph{mechanism} (where to inject), and the causal test in Step~3 remains decisive.
\end{resultbranch}

\paragraph{Step 2 --- steer, and freeze on dev (RQ6 operating point).} Steering adds the unit direction
at all generated positions,
\begin{equation}
h_\ell \leftarrow h_\ell+\alpha\,\vhat_\ell,\qquad \alpha=\rho\cdot\overline{\lVert h_\ell\rVert},
\label{eq:add}
\end{equation}
parameterized by the norm-relative $\rho$ so the coefficient transfers across layers (the ambient
residual norm grows with depth; a fixed $\rho$ injects a fixed fraction of it). Behavior is
\emph{non-monotone} in $\rho$---too much steering leaves the manifold and coherence collapses---so we
sweep $\rho$ under a coherence guardrail rather than fixing a large value. On a frozen 512-token eval
text we cap the mean per-token divergence of the steered from the unsteered next-token distribution,
$D_{\mathrm{KL}}(p_{\text{steered}}\|p_{\text{unsteered}})\le0.30$ nats, and require steered
render-success $\ge0.90$. The operating point is selected on \emph{dev only}:
\[
(\lstar,\rhostar,\text{variant}^{*})=\arg\max\ \text{dev POC-gain}\quad
\text{s.t.}\quad \text{mean-KL}\le0.30\ \text{and}\ \text{render-success}\ge0.90,
\]
ties broken toward lower KL. The dev sweep is Figure~\ref{fig:F6} (POC-gain over layer$\times\rho$,
guardrail-violating cells hatched) and the dose-response is Figure~\ref{fig:F7}. The frozen point is
$\lstar=\numslot{STEER-LSTAR}$, $\rhostar=\numslot{STEER-RHOSTAR}$, variant
$\numslot{STEER-VARIANT}$, with dev POC-gain $\numslot{STEER-DEVGAIN}$ at mean KL
$\numslot{STEER-KL}$ and render-success $\numslot{STEER-RENDER}$. It is written to a config file,
hashed, and git-tagged \emph{before} any held-out generation; the held-out set is touched exactly
once. If no single-layer config clears the guardrail-constrained threshold, we escalate to
multi-layer injection or a 2--4D subspace (pre-stated fallbacks) and re-select.

\paragraph{Step 3 --- the causal bar (RQ6), stated as interventions.} Steering and ablation are
$\doop$-operations on an activation node, severing its normal dependence on upstream computation.
The two arms are
\begin{align}
\textbf{sufficiency:}\quad&
\mathbb{E}\big[Q\mid \doop(h_{\lstar}\mathrel{+}{=}\alpha^{*}\vlstar),\ \NOSYS\big]\ >\ \mathbb{E}\big[Q\mid \NOSYS\big],
\label{eq:suff}\\
\textbf{necessity (flip):}\quad&
\mathbb{E}\big[Q\mid \doop(h\leftarrow Ph),\ \FULL\big]\ <\ \mathbb{E}\big[Q\mid \FULL\big],\qquad P=I-\vhat\vhat^{\top},
\label{eq:flip}
\end{align}
with $Q$ the POC or a preference utility. Sufficiency uses the \NOSYS{} base---\emph{no design
prompt}---so the effect is the direction's, not an instruction's; that is what makes it interventional
rather than observational. Necessity projects the direction out at every layer and position while the
skill is present (weight-orthogonalization, so no inference hook is needed), asking whether the skill
stops working. The projector $P$ is a genuine orthogonal projector (idempotent and symmetric;
Appendix~\ref{app:theory}). Held-out arms are unsteered-\NOSYS, steered-\NOSYS, a norm-matched random
control ($k{=}5$ draws re-normalized to $\lVert v_{\lstar}\rVert$), and an HF-generated \FULL{}
reference; we report both oracle signals and the reproduction fraction $\%\text{skill-reproduced}=
(\POC_{\text{steered}}-\POC_{\text{unsteered}})/(\POC_{\FULL}-\POC_{\text{unsteered}})\times100\%$.

\paragraph{The success criterion (verbatim).} The direction is a \emph{causally sufficient design
direction} iff, on held-out prompts (including 10 unseen-type):
\begin{quote}\itshape
steered-\NOSYS{} beats unsteered-\NOSYS{} on $\ge1$ objective family after Holm AND the
reliability-gated preference delta (steered $\succ$ unsteered) has BT-utility cluster-bootstrap CI
$>0$; AND the norm-matched random control ($k{=}5$ draws) does not satisfy both (specificity); AND
the flip test attenuates the skill's gain (necessity). If addition succeeds but the flip test does
not, report ``sufficient, not demonstrably necessary.''
\end{quote}
This four-conjunct standard is why a \emph{causal} headline is defensible: under a true null the joint
false-positive rate is below any single test's $\alpha$, and the four checks probe independent failure
modes---objective vs.\ preference vs.\ specificity vs.\ necessity (Appendix~\ref{app:theory}). The
specificity conjunct formalizes non-identifiability: if a random direction of equal norm also lifts
$Q$, the effect is a \emph{norm} effect, not a \emph{direction} effect. We additionally bound
side-effects: steering must not drop non-UI algorithmic pass-rate by $>10\%$ absolute nor push
general-text mean KL $>0.30$ nats.

\synthfig[0.47\linewidth]{F13}{F13_preview.pdf}{6}{Extraction stability: cosine of the running-mean
direction vs.\ number of examples, with the $\Delta\cos<0.01$ plateau line. Planted convergence;
the real curve certifies the 200-pair corpus (or triggers the seed-extension contingency).}

\synthfig[0.51\linewidth]{F5}{F5_preview.pdf}{6}{Where the skill signal lives: probe AUC, control-task
selectivity, and patch \%-recovered by layer, with the chosen band shaded and the AUC/selectivity
thresholds drawn. Planted mid-band plateau; the real curves set the Stage-2 focus layers (RQ5).}

\synthfig[0.47\linewidth]{F6}{F6_preview.pdf}{7}{Dev sweep: POC-gain over layer $\times\ \rho$;
guardrail-violating cells (KL$>0.30$ or render$<0.90$) hatched; the frozen $(\lstar,\rhostar)$ starred.
The operating point is selected here, on dev, then git-tagged before held-out. Planted optimum.}

\synthfig[0.51\linewidth]{F7}{F7_preview.pdf}{7}{Dose-response at $\lstar$: POC-gain (with CI band) and
render-success vs.\ $\rho$, failure region shaded. The non-monotone gain then coherence collapse is the
expected shape; planted here, measured on Colab.}

\section{Stage 2 results: does the direction steer?}
\label{sec:stage2results}

This is the pivotal section. With the operating point frozen on dev and git-tagged, we touch the
held-out set once and evaluate the four-conjunct criterion. The held-out arms---unsteered-\NOSYS,
steered-\NOSYS, norm-matched random ($k{=}5$), and an HF \FULL{} reference---are compared on both
oracle signals in Figure~\ref{fig:F8}; steering reproduces $\numslot{STEER-REPRO}$ of the skill's
POC gain. The result resolves into one of three preregistered branches, written out in full below.
Exactly one is kept on real data.

\subsection{Branch A --- a causally sufficient (and/or necessary) design direction exists}
\label{sec:branchA}

\begin{resultbranch}{A: causally sufficient direction}
Steering the diff-in-means direction with \emph{no design prompt} raises held-out UI quality on both
oracle signals: steered-\NOSYS{} beats unsteered-\NOSYS{} by $\numslot{STEER-DELTA}$ on the POC
(significant on $\numslot{STEER-NFAM}$ of six objective families after Holm) and is preferred with BT
utility $\numslot{STEER-BT}$ (cluster-bootstrap CI $>0$). Crucially, the norm-matched random control
does not clear both bars ($\numslot{STEER-RAND}$ POC, CI bracketing zero), so the effect is a
\emph{direction} effect, not a \emph{norm} effect---a random push of equal magnitude does not make
pages better. The direction reproduces $\numslot{STEER-REPRO}$ of the full skill's held-out gain
(Figure~\ref{fig:F8}), and it does so on prompts and \emph{task types the extraction corpus never
contained}. To our knowledge this is the first demonstration of a causal, steerable direction for the
visual/aesthetic quality of generated front-end code: a few sentences of design guidance and a
single mid-network vector turn out to move the same underlying quantity, and that quantity can be
actuated with the prose removed.

\smallskip\noindent\emph{Coherence and specificity.} The gain is not bought with gibberish: at
$\rhostar$ the mean next-token KL from the unsteered model is $\numslot{STEER-KL}$ nats (under the
$0.30$ guardrail) and steered render-success is $\numslot{STEER-RENDER}$. Steering leaves general
ability intact---non-UI algorithmic pass-rate drops by only $\numslot{SPEC-PASSDROP}$ (within the
$10\%$ tolerance) and general-text KL is $\numslot{SPEC-KL}$---so the direction is specific to design,
not a blunt instrument that degrades everything.
\end{resultbranch}

\paragraph{Necessity (the flip test).} Projecting $\vhat$ out at every layer and position while the
skill is present attenuates the skill's own gain by $\numslot{FLIP-ATTEN}$ (Eq.~\ref{eq:flip}), which
promotes the claim from sufficiency to necessity.
\begin{resultbranch}{A.1: sufficient and necessary}
The flip test attenuates the skill's gain by $\numslot{FLIP-ATTEN}$ with CI $>0$: removing the single
direction from the residual stream substantially disables the skill even when its 391 tokens are still
in context. The direction is therefore both sufficient (adding it with no prompt helps) and necessary
(removing it stops the prompt from helping)---the strong two-arm result, mirroring the refusal-direction
template but for a graded, multi-component aesthetic behavior. This is the paper's strongest mechanistic
claim: the skill's effect is, to a large extent, \emph{routed through this one direction}.
\end{resultbranch}
\begin{resultbranch}{A.2: sufficient, not demonstrably necessary}
Addition succeeds but the flip test does not attenuate the skill's gain ($\numslot{FLIP-ATTEN}$, CI
includes zero): the direction is a sufficient carrier but not the unique one. The skill routes its
effect through this direction \emph{and} others---consistent with a concept cone rather than a single
axis. We report ``sufficient, not demonstrably necessary'' and do not overclaim uniqueness; the honest
framing is ``\emph{a} causally sufficient design direction,'' never ``\emph{the} design direction.''
\end{resultbranch}

\paragraph{Failure surface and generalization (RQ7).} A positive mean can hide heavy-tailed per-prompt
behavior, so we report the distribution. The overall anti-steerable fraction (prompts whose steered POC
falls below unsteered) is $\numslot{FAIL-ANTI-OVERALL}$; on unseen types it is
$\numslot{FAIL-ANTI-UNSEEN}$, and the mean unseen-type steering delta is $\numslot{FAIL-UNSEEN-D}$
(Figure~\ref{fig:F12}). The reproduction fraction on unseen types is $\numslot{FAIL-REPRO-UNSEEN}$.
\begin{resultbranch}{RQ7 robust}
The anti-steerable fraction stays below $50\%$ across all task types, seen and unseen, and the
unseen-type delta is positive and gated: the direction transfers cleanly out-of-distribution. Steering a
direction extracted only from marketing- and content-oriented dev pages nonetheless improves
control-dense settings panels and admin tables, evidence that it encodes transferable design
\emph{content} rather than a task-idiom artifact.
\end{resultbranch}
\begin{resultbranch}{RQ7 partial}
Steering helps on average and generalizes to most unseen types, but a named regime is brittle:
$\numslot{FAIL-WORST-TYPE}$ crosses the $50\%$ anti-steerable line ($\numslot{FAIL-WORST-ANTI}$), so
steering makes those pages worse as often as better (Figure~\ref{fig:F12}). This is the honest failure
map---the direction is not a blanket win, and its data-dense failure mode is exactly where the aesthetic
prior (generous whitespace, one signature element) conflicts with the task's information density.
\end{resultbranch}
\begin{resultbranch}{RQ7 brittle}
The unseen-type effect collapses (anti-steerable $\ge50\%$ overall on OOD): the direction is
dev-specific and does not transport to genuinely new layouts. We scope the causal claim to the
in-distribution regime and report the OOD boundary precisely; the direction is real but local, and its
extraction corpus bounds where it works.
\end{resultbranch}

\paragraph{Component$\leftrightarrow$direction correspondence (RQ8).} Finally we test whether the
component structure the ablation found survives into the representation. Component sub-vectors
$v_{C_i}=\bar h_{\lstar}^{(\AOI\text{-}C_i)}-\bar h_{\lstar}^{(\NEUTRAL)}$ have mean off-diagonal
cosine $\numslot{CORR-COSMEAN}$ (max $\numslot{CORR-COSMAX}$), and $\numslot{CORR-NMATCH}$ of five steer
their own metric family (Figure~\ref{fig:F9}). Cosine can show sub-vectors are linearly distinguishable
but not that each \emph{causally} steers its family, so the rule pairs near-orthogonality with a
steering-signature match.
\begin{resultbranch}{RQ8 found}
The sub-vectors are near-orthogonal (mean $|\!\cos|=\numslot{CORR-COSMEAN}<0.3$) \emph{and}
$\ge3/5$ steer their own metric family---C1 moves color, C2 moves layout, and so on. The design skill
has a \emph{component basis} inside the model, and the black-box ablation signatures correspond to
white-box steering directions: the paper's strongest bridge result, an ablation$\leftrightarrow$activation
correspondence predicted by the geometry of linear representations. It means ``clean design'' is not one
entangled thing but a small set of separable, individually actuatable sub-concepts.
\end{resultbranch}
\begin{resultbranch}{RQ8 partial}
One or two sub-vectors steer their own family while the rest are entangled ($\numslot{CORR-NMATCH}/5$
match): the representation has partial component structure. Some components (likely color and layout,
the most necessary) have distinct directions; others share a subspace. The bridge holds in part---the
most causal components are the most separable---but ``clean design'' is not a fully diagonal basis.
\end{resultbranch}
\begin{resultbranch}{RQ8 absent}
The sub-vectors are not near-orthogonal, or fewer than three steer their own family: there is no
component basis. ``Clean design'' is one entangled direction inside the model even though the prompt-level
skill decomposes into five components. This is still informative and slightly surprising: the ablation's
component structure is a property of the \emph{instruction}, not of the model's internal
representation---the model compresses five instructions into one axis of variation.
\end{resultbranch}

\subsection{Branch B --- the honest null}
\label{sec:branchB}

\begin{resultbranch}{B: not linearly steerable}
No single linear direction steers held-out UI quality under the coherence guardrail. Steered-\NOSYS{}
does not beat unsteered-\NOSYS{} on any objective family after Holm ($\numslot{STEER-DELTA}$ POC, CI
$\numslot{STEER-DELTA-CI}$ bracketing zero) and the preference delta's BT CI includes zero
(Figure~\ref{fig:F8}). This is a \emph{null with teeth}, not an absence of evidence: the power gate
certifies the confidence intervals are tight enough to exclude effects the study was designed to
detect, and the extraction direction \emph{is} decodable (RQ5 probes separate \FULL{} from \NEUTRAL{}
at high AUC)---so we have a direction that \emph{correlates} with design but does not \emph{cause} it
when injected. A correlational direction that does not steer is a hypothesis, not a finding; we report
it as such.

\smallskip\noindent The evidential value is precise. The null \emph{rules out} the strongest form of
the linear-representation account for this behavior in this model: that a single mid-network direction,
added with no prompt, suffices to reproduce the skill's effect. It is consistent with two readings we
cannot distinguish here and state honestly: either the skill's effect is \emph{distributed} across
many directions and/or layers (a single injection is too blunt---the pre-stated multi-layer and
subspace fallbacks were tried and also failed), or the behavior is genuinely \emph{non-linear} in the
residual stream (diff-in-means, a linear probe of a nonlinear model, cannot capture it). Either way,
the practical implication is that the design skill's mechanism is not a simple steerable feature: the
391 tokens do more than trigger one latent switch. For the interpretability literature, it is a
negative datapoint on how far the single-direction picture---so clean for refusal---extends to graded,
compositional, generation-spanning aesthetic behaviors. We publish it with the same rigor as a positive
result, because a forced or buried null would corrupt exactly the oracle this project is built to
defend.
\end{resultbranch}

\subsection{Branch C --- partial outcomes}
\label{sec:branchC}
Between a clean causal direction and a clean null lie four partial results, each shipped with its
numbers and each carrying a distinct, bounded claim.

\begin{resultbranch}{C.1: steers the objective signal only}
Steered-\NOSYS{} beats unsteered on $\ge1$ objective family after Holm ($\numslot{STEER-DELTA}$), but
the preference delta's BT CI includes zero. The direction moves the \emph{measurable} correlates of
design---contrast, alignment, slop tells---without moving \emph{human-judged} quality. Given that
objective metrics cap at $\sim$half the variance in appeal, this is a meaningful boundary: we have
steered what the metrics see but not what a designer prefers, and by the null rule the aesthetic claim
is inadmissible. We report a ``metric-level'' steering effect and explicitly decline the quality claim.
\end{resultbranch}
\begin{resultbranch}{C.2: steers, but only by breaking coherence}
A configuration raises POC only above the KL guardrail ($\numslot{STEER-KL}$ nats $>0.30$) or below the
render floor: the direction ``improves'' pages by pushing the model off-distribution, and the gain
rides on artifacts rather than design. Under the preregistered guardrail this does not count as
success. We report the direction as \emph{coherence-breaking} and show the dose-response
(Figure~\ref{fig:F7}) where the apparent gain and the KL blow-up coincide---an honest account of a
non-result that a laxer guardrail would have mis-sold as a finding.
\end{resultbranch}
\begin{resultbranch}{C.3: anti-steerable majority}
The mean steering effect is near zero but hides a heavy-tailed distribution: $\numslot{FAIL-ANTI-OVERALL}$
of prompts are anti-steerable, so steering helps half and harms half (Figure~\ref{fig:F12}). The
``average effect'' is an artifact of cancellation, and reporting only the mean would be actively
misleading---the reason the protocol reports the full per-prompt distribution. The direction is not a
reliable design lever; whatever it encodes is prompt-contingent in a way a single global vector cannot
resolve.
\end{resultbranch}
\begin{resultbranch}{C.4: sufficient but not necessary}
Addition succeeds on both signals (Branch~A holds) but the flip test does not attenuate the skill's
gain ($\numslot{FLIP-ATTEN}$, CI includes zero): a steerable direction exists, yet the skill does not
\emph{depend} on it---the model has other routes to the same behavior. This is the concept-cone
outcome; we claim a sufficient direction and an honest account of non-uniqueness, and we present the
cosine/cone structure (Figure~\ref{fig:F9}) rather than a single-axis story.
\end{resultbranch}

\synthfig[0.47\linewidth]{F8}{F8_preview.pdf}{8}{The held-out causal test. POC by arm---unsteered,
steered (no design prompt), norm-matched random ($k{=}5$), and \FULL{} reference---with the
\%-skill-reproduced annotation. Success requires steered $>$ unsteered while random $\approx$ unsteered.
Planted arms illustrate the layout; the real bars decide RQ6.}

\synthfig[0.53\linewidth]{F12}{F12_preview.pdf}{8}{Anti-steerable fraction by task type (seen vs.\
unseen/OOD), with the $50\%$ brittleness line. A positive mean can coexist with a brittle unseen type;
the distribution, not the mean, is the RQ7 claim. Planted; real data replaces the bars.}

\synthfig[0.7\linewidth]{F9}{F9_preview.pdf}{9}{Component$\leftrightarrow$direction correspondence
(RQ8). Left: cosine matrix among the five component sub-vectors and $v_{\text{full}}$ (near-zero
off-diagonal $\Rightarrow$ orthogonal). Right: steer-signature $\times$ ablation-signature match (a
bright diagonal $\Rightarrow$ each component steers its own family). Planted diagonal; real data
decides found/partial/absent.}

\section{Discussion}
\label{sec:discussion}

\paragraph{The bridge: does the black-box explanation match the white-box one?} The study's
distinctive move is to ask one question two ways---which \emph{instruction} matters (Stage~1) and
which \emph{direction} matters (Stage~2)---on the same model, and then to check whether the answers
agree. That check is the correspondence between the components RQ1 finds most necessary and the
sub-directions RQ8 finds separable.
\begin{resultbranch}{bridge holds}
If RQ1 ranks C5 (negatives) and C1 (color) most necessary \emph{and} RQ8 finds those the sub-vectors
that most cleanly steer their own metric families, the two explanations converge: the parts of the
prose that carry the causal weight are the parts with the crispest internal directions. That
convergence is the paper's most satisfying result---a prompt-level account and a
representation-level account of the \emph{same} phenomenon, mutually corroborating---and it suggests
skill authors are, in effect, writing instructions that map onto the model's existing axes of
variation.
\end{resultbranch}
\begin{resultbranch}{bridge fails}
If the most necessary components are \emph{not} the most separable directions (RQ8 absent or
mismatched), the two levels tell different stories: the instruction's compositional structure is not
the model's. This is equally interesting---it says the componentization that helps a human author a
skill is not how the model represents design---and it cautions against reading prompt structure as
mechanism.
\end{resultbranch}

\paragraph{What ``make it beautiful'' tells us about instruction compression.} RQ4 turns a cheap
one-liner into a probe of how much of a 391-token skill is \emph{knowledge} versus \emph{permission}.
If a five-word nudge recovers most of the effect (RQ4 nudge), the model already contains the
capability and the skill largely licenses its use---which is precisely the premise that makes Stage~2
plausible: a capability the model already has is one a direction can actuate. If instead the
structured content is load-bearing (RQ4 content), the skill transfers specific
knowledge---named palettes, type scales, prohibitions---and we should expect the extracted direction
to be correspondingly richer than a single ``try harder'' axis. Either way, RQ4 and Stage~2 are two
readings of the same question: is the design skill flipping a switch the model already has, or
installing something new?

\paragraph{Mechanistic implications.} A positive Stage~2 result would extend the single-direction
picture---clean for refusal \citep{arditi2024refusal}---to a graded, compositional, generation-spanning
\emph{aesthetic} behavior, which is a substantially harder case than a binary safety behavior: the
target is downstream of thousands of tokens of code and has no single ``answer token.'' A null would
mark a boundary of that picture, and the honest reading is that not every behavior a prompt induces is
a steerable feature. The correspondence result (RQ8), whichever way it falls, speaks to the geometry of
compositional concepts \citep{park2024geometry}: whether a human's five-way decomposition of ``design''
lines up with the model's internal near-orthogonal directions.

\paragraph{Relation to the steering-reliability literature.} We inherit its central warning: mean
effects hide heavy tails, and up to half of inputs can be anti-steerable
\citep{tan2024analysing,dasilva2025steeringoff}. Our protocol answers it directly---per-prompt
distributions, an anti-steerable fraction, a class-separation pre-screen, an OOD held-out over unseen
\emph{types}, and a coherence guardrail on a non-monotone dose-response
\citep{taimeskhanov2026strength}. The four-conjunct success criterion, in particular, is designed so
that a strong average cannot by itself buy a causal claim: specificity and necessity must also hold.
This is the methodological transfer we hope outlasts the particular numbers---a template for claiming a
\emph{causal} steering result responsibly when the underlying phenomenon is known to be unreliable.

\paragraph{Practical guidance for skill design.} Read together, the Stage-1 outcomes yield actionable
advice for authors of design skills: invest tokens where necessity concentrates (the RQ1 ranking),
prune any component whose leave-one-out is null or negative, and exploit the interactions RQ3
surfaces---if C1$\times$C5 is super-additive, a palette instruction and a ``no purple gradients''
prohibition should ship together, not separately. If sufficiency is distributed (RQ2), resist the urge
to shorten the skill to its single largest component; if it is concentrated, the skill can be
compressed. These are empirical recommendations, contingent on which branch obtains, and each is
traceable to a preregistered contrast.

\section{Limitations and threats to validity}
\label{sec:limitations}

We state the bounds of the study up front, and pair each internal-validity threat with the control
that neutralizes it.

\paragraph{One model, one family.} Every result is about \texttt{Qwen2.5-Coder-7B-Instruct}. The
two-stage design \emph{requires} a single open-weights model so that the ablation and the
interpretability describe the same system; the price is that generalization to other models and
scales is untested. This is a deliberate scope choice, not an oversight, and the recorded fallback
model would restart both stages rather than mix them.

\paragraph{One skill instantiation.} We ablate one five-component skill with one content-orthogonal
decomposition (C1--C4 prescriptive, C5 all-negatives). The componentization is one defensible cut
among several---a different author might split typography from type-scale, or fold patterns into
layout---and the necessity/sufficiency verdicts are conditional on this cut. We fix and document the
decomposition and its token balance, and note that RQ8 tests whether \emph{this} cut corresponds to
internal structure; a different cut could correspond differently.

\paragraph{The judge is not a human panel.} The preference signal is a VLM judge validated on a
100-pair human subset and gated on chance-corrected agreement, but it is not a full human study. Its
known biases (length, verbosity, position) are controlled by both-orders judging and style-controlled
Bradley--Terry, and the objective half is the unfoolable anchor---but the preference channel's ceiling
is the judge's fidelity, which the reliability gate bounds rather than eliminates. Where the two
$\hat\beta$ (controlled vs.\ uncontrolled) disagree in sign, we report the contrast as style-confounded
rather than resolved.

\paragraph{Objective metrics cap at half the variance.} Interface-aesthetics metrics explain
$\lesssim$half the variance in human appeal \citep{miniukovich2015computation,reinecke2013predicting}.
We therefore never let the objective half stand alone for an aesthetic claim; it is necessary
(deterministic, unfoolable) but not sufficient, which is the whole reason for the second signal. The
PSI is validated against human ``AI-slop'' ratings before use, and demoted to a 6-term POC if it fails.

\paragraph{Diff-in-means is a linear probe of a nonlinear model.} A null steering result bounds
\emph{linear} steerability at the chosen sites, not the existence of the concept: the behavior could be
non-linear \citep{engels2024notlinear} or distributed across a subspace or cone
\citep{wollschlager2025cones}. We say so explicitly, try the multi-layer and subspace fallbacks before
declaring a null, and frame any positive result as ``\emph{a} sufficient direction,'' never ``the''
direction---non-identifiability is assumed, and the specificity control ($k{=}5$ random draws) is what
separates a direction effect from a norm effect.

\paragraph{The interpretability threats, and their controls.} Table~\ref{tab:threats} lists the
Stage-2 threats and the control that neutralizes each; these are built into the design, not applied
afterward.

\begin{table}[htbp]
\centering\footnotesize
\caption{Stage-2 internal-validity threats and controls (Table~T9, abridged).}
\label{tab:threats}
\begin{tabular}{@{}p{0.30\linewidth}p{0.62\linewidth}@{}}
\toprule
threat & control \\
\midrule
Norm confound (any large push jogs quality) & norm-matched random vector, $k{=}5$ draws; must not reproduce the gain \\
Non-identifiability / concept cones & claim ``a'' not ``the'' direction; report cosine/cone structure (RQ8) \\
Self-repair inflating patching necessity & \emph{denoising} (not noising) activation patching \\
Interpretability illusion (dormant pathway) & triangulate probing + patching + steering (RQ5 rule) \\
Spurious extraction cues & held-out incl.\ unseen task \emph{types}; class-separation pre-screen \\
Selection leakage (tuning on test) & freeze $(\lstar,\rhostar,\text{variant})$ on dev, git-tag before held-out, touch once \\
Engine confound (vLLM vs.\ HF) & within-HF comparisons only; re-generate baselines in HF \\
\bottomrule
\end{tabular}
\end{table}

\paragraph{Statistical scope.} Cluster-bootstrap CIs are calibrated to the \emph{number of prompts}
(40 dev, 18 held-out), not the number of generations, so their tails are limited at 40 clusters; we
report parametric MixedLM CIs as a robustness pair and do not claim generalization beyond the 58-brief
task population. Post-hoc power is used only as an admissibility screen, never to reinterpret an
under-powered delta as a null.

\paragraph{Compute scale.} The GPU budget is modest (a few L4- and A100-hours; Section~\ref{sec:repro}),
which fixes the corpus and seed counts. Larger corpora or more seeds would tighten every CI; the design
allocates its five headline seeds and 40 dev prompts to the confirmatory contrasts and accepts lighter
power on the exploratory scans, controlled by BH.

\section{Reproducibility and ethics}
\label{sec:repro}

\paragraph{Preregistration and freezes.} All confirmatory decisions---the cell table, the 58-brief
split, the POC formula and its two conditionals, the 12-test confirmatory family, the reliability and
power gates, the Stage-2 four-conjunct success criterion, and the null rule---are frozen in a
preregistration (Appendix~\ref{app:prereg}) git-tagged \texttt{prereg-v1} before the first GPU run.
The Stage-2 operating point $(\lstar,\rhostar,\text{variant}^{*})$ is selected on dev, written to a
hashed config, and git-tagged \texttt{steer-frozen} before the held-out set is touched; held-out is
evaluated exactly once. Any post-freeze change is a dated amendment; a confirmatory change forks a new
tag and is reported as a deviation.

\paragraph{Determinism and pinning.} An engine split is enforced: vLLM for Stage-1 bulk generation, HF
for all activation and steering work, never cross-compared; engine, version, sampling, and config
hashes are stamped on every results row. Sampling is fixed (Section~\ref{sec:design}). Rendering is
deterministic: pinned Chromium at $1440\times900$ (plus a $390\times844$ mobile subset for the
responsiveness slice), device-scale-factor 1, animations disabled, fonts-ready awaited,
network blocked; a metric artifact re-run under an unchanged config reproduces identical hashes. All
pinned versions are recorded in an environment lock.

\paragraph{Artifacts and machinery.} Reusable machinery lives in a tested Python package (metrics,
rendering, judging, agreement, Bradley--Terry, mixed models, power, hooks, steering, figures,
manifest); notebooks carry narrative and orchestration and import from it. The CPU-safe half---fixture
rendering, the full objective-metric stack, the judge in dry-run/mock mode, the statistics on
synthetic data, and every unit test---runs and passes in the development environment; the GPU-dependent
half is written, mock-tested, import-guarded, and left unexecuted until the Colab sessions in the
runbook. Every figure in this paper is produced by the same figure module from a results table---no
hand-editing---and the preview figures here are that module run on synthetic data with a planted ground
truth (unit-tested to recover the plant). An artifact manifest records every HTML, screenshot,
DOM-JSON, activation shard, and results file with a SHA-256, its producing config hash, engine, and
commit.

\paragraph{Cost.} The full study is deliberately inexpensive (Table~\ref{tab:cost}), which we regard
as a feature: a two-stage causal study of a design skill fits in a graduate compute budget.

\begin{table}[htbp]
\centering\footnotesize
\caption{Compute and API budget (Table~T7). These are planned budgets, not results.}
\label{tab:cost}
\begin{tabular}{@{}lll@{}}
\toprule
item & quantity & cost \\
\midrule
Stage-1 generation (vLLM, L4) & $\sim$6 L4-hours & \multirow{2}{*}{$\approx$243 Colab units, \$50--65} \\
Stage-2 capture/sweep/verify/stretch (A100) & $\sim$14 A100-hours & \\
Judge --- Gemini 2.5 Flash primary & $\sim$4{,}420 both-order calls & $\approx$\$7 \\
Judge --- GPT-4o audit (15\%) & $\sim$664 calls & $\approx$\$7 \\
Storage (HTML/PNG/DOM-JSON/activations) & $\sim$6{,}900 units + means & $\approx$4--5 GB \\
\midrule
\textbf{Total} & & \textbf{$\approx$\$65--80} \\
\bottomrule
\end{tabular}
\end{table}

\paragraph{Ethics.} The study generates web-UI code and evaluates its visual quality; it involves no
personal data and no human subjects beyond the authors' own labeling of a 100-pair validation subset
(no demographic data, no external participants). The subject model is used under its license; API
judges are used only as evaluators. The one dual-use-adjacent technique is weight-orthogonalization,
used here solely as the necessity flip test on an \emph{aesthetic} direction, with no safety-relevant
behavior involved. We see no additional ethical concerns.


\small
\bibliography{refs}

\normalsize
\appendix
% ============================ APPENDICES ============================

\section{The design skill, fillers, and padding pool (verbatim)}
\label{app:skill}

The frozen independent variable. Component boundary tags \texttt{[[Cn]]} are documentation only and
are stripped before the text enters the system prompt; canonical order is C1$\to$C2$\to$C3$\to$C4$\to$C5.
C1--C4 are purely prescriptive; C5 consolidates all negative constraints.

\paragraph{C1 --- color / palette.}
\begin{quote}\small
Choose a deliberate palette of four to six specific colors --- one dominant hue, one or two supporting
tones, and one sharp accent --- each fixed as an exact hex value and reused consistently throughout.
Build depth with atmospheric, layered backgrounds such as subtle tints, soft shading, or fine texture
rather than flat fills, and ensure every text-on-background pairing clears WCAG AA contrast.
\end{quote}
\paragraph{C2 --- layout / spacing / hierarchy.}
\begin{quote}\small
Open with the most characteristic element of the subject so the page announces what it is at a glance.
Establish a clear visual hierarchy on a consistent spacing scale, aligning elements to a small set of
shared edges so the composition reads as engineered. Favor intentional, sometimes asymmetric
arrangements over uniform grids, and give the layout generous, purposeful white space so it can breathe.
\end{quote}
\paragraph{C3 --- typography.}
\begin{quote}\small
Pair a characterful display typeface, used sparingly for headings, with a clean, highly readable body
typeface and a compact utility face for small labels and data. Set a clear modular type scale with a few
deliberate sizes and weights, and let the typography carry real personality --- expressive at large
sizes, quietly legible at small ones --- instead of neutrally delivering text.
\end{quote}
\paragraph{C4 --- component patterns.}
\begin{quote}\small
Design one well-chosen signature element --- an unexpected section, a distinctive card treatment, or a
considered navigation --- that anchors the page's identity. Add at most one orchestrated motion moment,
such as a single page-load reveal or a refined hover state, and make it feel deliberate. Vary corner
radii, borders, and shadows with intent so that depth and emphasis track the importance of the content.
\end{quote}
\paragraph{C5 --- negative constraints (all negatives).}
\begin{quote}\small
Avoid AI slop defaults: purple or indigo gradients; Inter and Roboto typefaces; fully centered hero
layouts built around one big headline and stat; three identical icon cards in a row; uniform rounded
corners on everything; and scattered or excessive animation. Never default to generic template palettes
or cream with terracotta serif and near black with acid green looks, unless the brief calls for them.
Cut any decoration that serves no purpose.
\end{quote}

\paragraph{Neutral fillers (render-inert; matched to C1--C5).} A filler may touch only source-level
conventions and is forbidden from mandating anything an objective metric could detect (semantic
landmarks, \texttt{lang}/\texttt{title}/\texttt{meta}, ARIA, validation, event handlers), enforced by a
build-time audit against a banned-topic list.
\begin{quote}\small
\textbf{Filler-1} ($\leftrightarrow$C1): Write the markup in a consistent source style: use lowercase
for tag and attribute names, quote every attribute value with double quotes, and keep a stable attribute
order within each tag. Break long lines at a consistent indent so no line grows unwieldy, and keep
quoting and casing choices uniform across the whole file so the source reads cleanly.\\[2pt]
\textbf{Filler-2} ($\leftrightarrow$C2): Indent the markup consistently with two spaces per level and
keep lines to a reasonable width so the source stays legible. Group related rules together, order the CSS
declarations predictably, and leave a short comment above each major part of the stylesheet describing
what it styles. Avoid trailing whitespace and keep a single blank line between logical blocks.\\[2pt]
\textbf{Filler-3} ($\leftrightarrow$C3): Name classes and identifiers with clear, lowercase,
hyphen-separated words that describe purpose rather than appearance, and keep the naming scheme
consistent across the file. Avoid abbreviations that are not obvious, do not reuse one name for two
purposes, and prefer a short, stable vocabulary of names over inventing a new term for every element on
the page.\\[2pt]
\textbf{Filler-4} ($\leftrightarrow$C4): Organize any JavaScript into small, single-purpose functions
with descriptive names, each declared before the point where it is used, and keep the logic shallow
rather than deeply nested. Prefer const and let over var, keep variable declarations near their first
use, group related functions together in the script, and leave a brief comment explaining any step whose
intent is not obvious.\\[2pt]
\textbf{Filler-5} ($\leftrightarrow$C5): Comment the code where intent is not obvious, but avoid
restating what the code plainly does; a good comment explains why, not what. Keep comments short, place
each one directly above the line or block it describes, and use a consistent comment style throughout.
Remove commented-out fragments and leftover notes to self before finishing, so only purposeful comments
remain in the final file.
\end{quote}

\paragraph{Padding pool (frozen; build-time equalization).} Each filler is deterministically padded (in
fixed order, skipping any clause that would overshoot) or trimmed at a sentence boundary until it matches
its component's Qwen token count to within $\pm2$: \emph{Pad-1a} Apply the same wrapping style to every
long tag so no line stands out; \emph{1b} Break attribute lists at a consistent width; \emph{1c} Keep
tag casing uniform throughout; \emph{2a} Keep spacing between rule blocks even so the stylesheet reads in
a steady rhythm; \emph{2b} Order declarations the same way in every rule; \emph{2c} Keep rule blocks tidy
and short; \emph{3a} Prefer plain descriptive words over clever coinages when a name must be introduced;
\emph{3b} Keep names short, plain, and easy to scan; \emph{3c} Keep the vocabulary of names small;
\emph{4a} Group constants near the top of the script so their values are easy to locate; \emph{4b} Keep
the script's functions in one predictable order; \emph{4c} Prefer flat logic over nesting; \emph{5a}
Delete stale comments rather than letting them drift out of date; \emph{5b} Keep comment punctuation
simple and consistent; \emph{5c} Prefer one clear comment over three vague ones.

\paragraph{Constant output constraint (in the user turn; never ablated).}
\begin{quote}\small
Return a single self-contained HTML file: put all CSS in one \texttt{<style>} tag and all JavaScript in
one \texttt{<script>} tag, inline, with no external network requests --- no external fonts, stylesheets,
scripts, images, or CDNs. Use only system or generic fonts, and inline SVG or CSS for any graphics.
Output only the HTML file.
\end{quote}

\paragraph{Token balance.} Exact Qwen2.5-Coder tokens per component: C1 75, C2 73, C3 75, C4 80, C5 88
(mean $m{=}78.2$, $\pm15\%$ band $[66.5,89.9]$; all within band); full skill $=$ 391 exact tokens.
Equalized fillers: F1 75, F2 74, F3 76, F4 81, F5 87 (each within 2 tokens of its component). The four
control cells are \FULL{} (C1--C5), \NEUTRAL{} (Filler-1--5), \BEAUTY{} (``Make it beautiful and
well-designed.''), and \NOSYS{} (no system message).

% ------------------------------------------------------------------
\section{Prompt corpus (frozen by id)}
\label{app:corpus}

The 58 briefs (Table~T2). Difficulty E/M/H by element count and interaction complexity. Split:
\texttt{dev} (40), \texttt{seen} (held-out, seen type; 8), \texttt{unseen} (held-out, OOD type; 10).
All briefs implicitly carry the constant output constraint; a whole-word leakage audit over the briefs
passes (zero matches of the banned style/design word list).

\begin{small}
\begin{longtable}{@{}p{0.045\linewidth}p{0.075\linewidth}p{0.70\linewidth}p{0.02\linewidth}p{0.055\linewidth}@{}}
\toprule
id & type & brief & d & split \\
\midrule
\endfirsthead
\toprule id & type & brief & d & split \\ \midrule \endhead
\bottomrule \endfoot
L01 & landing & Landing page for a mobile app that tracks home water usage and flags leaks; headline, how-it-works, three feature highlights, download call. & E & dev \\
L02 & landing & Landing page for a regional apple-cider producer selling seasonal boxes; product story, what's in a box, delivery areas, order button. & M & dev \\
L03 & landing & Landing page for a B2B API that turns shipping documents into structured data; hero statement, integration steps, a code sample, request-access form. & H & dev \\
L04 & landing & Launch page for an independent documentary about deep-sea mining; synopsis, trailer placeholder, screening dates, newsletter signup. & H & dev \\
L05 & landing & Landing page for a neighborhood bicycle-repair co-op; hours, services, a pricing note, location. & E & dev \\
L06 & landing & Landing page for a language-exchange meetup pairing learners for weekly conversation; how pairing works, supported languages, join button. & M & seen \\
D01 & dashboard & Dashboard for a solar-panel owner; today's generation, home consumption, battery level, 7-day history chart placeholder. & M & dev \\
D02 & dashboard & Operations dashboard for a bike-share system; station occupancy, bikes in transit, low-battery e-bikes, alerts list. & H & dev \\
D03 & dashboard & Personal-finance dashboard; account balances, monthly spending by category, upcoming bills, savings-goal tracker. & M & dev \\
D04 & dashboard & Dashboard for a small warehouse; inventory by aisle, orders awaiting pick, staff on shift, throughput chart placeholder. & H & dev \\
D05 & dashboard & Dashboard for a habit-tracking app; streaks for four habits, a weekly completion grid, today's checklist. & E & dev \\
D06 & dashboard & Fleet dashboard for a delivery company; vehicles active, deliveries completed, fuel usage, map placeholder with status pins. & M & seen \\
F01 & form & Multi-step signup form for a coding bootcamp; personal details, program choice, payment plan, progress indicator. & E & dev \\
F02 & form & Grant-application form for a community arts fund; applicant info, project description, budget table, file-upload fields. & M & dev \\
F03 & form & Booking form for a pottery studio; class selection, date and time, seats, contact details. & M & dev \\
F04 & form & Renters-insurance quote form; address, coverage options, valuables to insure, a summary step before submitting. & H & dev \\
F05 & form & Contact form for a wedding photographer; names, event date, venue, package interest, message. & E & dev \\
F06 & form & Patient-intake form for a dental clinic; personal and insurance details, medical-history checkboxes, consent. & M & seen \\
P01 & pricing & Pricing page for a note-taking app; three tiers (free/personal/team), monthly-annual toggle, feature comparison. & E & dev \\
P02 & pricing & Pricing page for a car-wash membership; three plans, what each includes, an FAQ section. & M & dev \\
P03 & pricing & Pricing page for a freelance illustrator; service packages, add-ons, turnaround times, booking link. & M & dev \\
P04 & pricing & Pricing page for a cloud-storage service; usage-based tiers, an interactive cost estimate, enterprise contact option. & H & dev \\
P05 & pricing & Pricing page for a yoga studio; drop-in, class packs, unlimited monthly, a short note on each. & E & dev \\
P06 & pricing & Pricing page for a project-management tool; four plans across a long feature matrix with category groupings. & M & seen \\
PF01 & portfolio & Portfolio home for a furniture maker; project grid, short bio, contact details. & E & dev \\
PF02 & portfolio & Portfolio for a film-score composer; reel placeholder, selected credits, a list of instruments and tools used. & M & dev \\
PF03 & portfolio & Portfolio for an architecture studio; three built projects with descriptions, a studio statement. & M & dev \\
PF04 & portfolio & Portfolio for a data journalist; long-form interactive pieces, each with a synopsis and a role note. & H & dev \\
PF05 & portfolio & Portfolio for a tattoo artist; a gallery, studio location, a booking enquiry link. & E & dev \\
PF06 & portfolio & Portfolio for a motion-graphics freelancer; work organized into categories, a short about section. & M & seen \\
B01 & blog & Article page about restoring a vintage typewriter; headings, image placeholders, an author note. & E & dev \\
B02 & blog & Blog index for a gardening site; recent articles by season, tags, a search field. & M & dev \\
B03 & blog & Long-form article on the history of the tea trade; table of contents, pull quotes, footnotes. & M & dev \\
B04 & blog & Magazine-style feature profiling a wildlife photographer; full-width imagery placeholders mixed with narrative columns. & H & dev \\
B05 & blog & Blog post reviewing three noise-cancelling headphones; a verdict box, pros and cons. & E & dev \\
B06 & blog & Recipe article for a sourdough loaf; ingredients, step-by-step method, timing notes. & M & seen \\
E01 & ecommerce & Product page for a stainless-steel water bottle; image gallery, price, size and color options, add-to-cart. & E & dev \\
E02 & ecommerce & Product page for a mechanical keyboard; switch options, specifications, reviews summary, related products. & M & dev \\
E03 & ecommerce & Product listing for a bookstore's staff picks; filters by genre, sort options, a grid of titles. & M & dev \\
E04 & ecommerce & Configurator page for a made-to-order desk; material, dimensions, cable options, live price, order. & H & dev \\
E05 & ecommerce & Product page for single-origin coffee; roast details, tasting notes, grind selection, subscribe option. & E & dev \\
E06 & ecommerce & Product listing for a plant nursery; filtered by light and care level, each card showing name and price. & M & seen \\
A01 & auth & Login screen for a budgeting app; email and password, remember-me, reset-password link. & E & dev \\
A02 & auth & Registration screen for a co-working booking site; email, password-strength meter, terms acceptance. & M & dev \\
A03 & auth & Two-factor screen prompting for a six-digit code; resend option, a help note. & M & dev \\
A04 & auth & Account-recovery flow screen; multiple verification methods, explanation of what happens next. & H & dev \\
A05 & auth & Sign-in for a library catalog; card number and PIN, a guest-access option. & E & dev \\
A06 & auth & Single-sign-on chooser screen; pick an organization before continuing to log in. & M & seen \\
ST01 & settings & Settings page for a chat app; notifications, privacy, appearance, connected devices, using labeled toggles and selects. & M & unseen \\
ST02 & settings & Account settings for a streaming service; profile, subscription, playback quality, parental controls, downloads. & H & unseen \\
ST03 & settings & Workspace settings for a team tool; members, roles and permissions, billing, integrations. & M & unseen \\
ST04 & settings & Notification-preferences page; choose channels (email/push/SMS) per category, a save bar. & E & unseen \\
ST05 & settings & Device settings for a smart thermostat; schedules, temperature units, eco mode, sensor calibration. & H & unseen \\
AT01 & admin-table & Admin users table; name, email, role, status, last-active columns, search, row actions. & M & unseen \\
AT02 & admin-table & Admin orders table; filterable columns, status badges, bulk selection, pagination, detail-drawer trigger. & H & unseen \\
AT03 & admin-table & Admin content-moderation queue; reported items with reason, reporter, date, approve/remove actions. & M & unseen \\
AT04 & admin-table & Admin inventory table; SKU, product, stock, price, an edit action per row, a filter bar. & E & unseen \\
AT05 & admin-table & Admin analytics table; campaigns with sortable metric columns, sparkline placeholders, an export control. & H & unseen \\
\end{longtable}
\end{small}

% ------------------------------------------------------------------
\section{Statistical and interpretability derivations}
\label{app:theory}

Notation is identical to the project's theory appendix; each result names the section it lifts.

\paragraph{C.1 Factorial model, and the necessity/sufficiency identity (T1).} Code a cell by
$x\in\{\pm1\}^5$. The effects parameterization of the $2^5$ writes the expected response as the
multilinear polynomial $\mathbb{E}[y(x)]=\beta_0+\sum_i\beta_i x_i+\sum_{i<j}\beta_{ij}x_ix_j+\cdots$.
Evaluating at the relevant corners (only $x_i$ flips; partners held at $+1$ for necessity, $-1$ for
sufficiency) gives Eq.~\ref{eq:necsuff}: the main term contributes $2\beta_i$ in both, while a pair
$\{i,j\}$ contributes $+2\beta_{ij}$ when the partner is present and $-2\beta_{ij}$ when absent. Hence
$\tfrac12(\mathrm{Nec}_i+\mathrm{Suff}_i)=2\beta_i+R$ and
$\tfrac14(\mathrm{Nec}_i-\mathrm{Suff}_i)=\sum_{j\neq i}\beta_{ij}$ (exact through three-factor terms,
which are equal in both contrasts and cancel in the difference). Necessity and sufficiency thus coincide
iff a component's interactions vanish; synergy makes a component more necessary than sufficient.

\paragraph{C.2 Resolution-V aliasing (T2).} The half-fraction with generator $E{=}ABCD$ has defining
subgroup $\{I,ABCDE\}$; the alias of an effect $S$ is $S\cdot ABCDE$ reduced modulo $x_i^2{=}1$. The
shortest defining word has length 5, so the design is resolution~V: main effects alias only with
four-factor interactions ($A{=}BCDE$) and two-factor interactions only with three-factor interactions
($AB{=}CDE$). Under ``three-factor-and-higher negligible,'' all five mains and all ten 2FIs are cleanly
and simultaneously estimable from 16 runs.

\paragraph{C.3 Mixed model and pseudoreplication (T3).} With $y_{cps}=\mu+\tau_c+u_p+v_{cp}+
\epsilon_{cps}$, two seeds on a prompt share $u_p$, so the intraclass correlation is
$\rho_{\mathrm{ICC}}=(\sigma_p^2+\sigma_{cp}^2)/(\sigma_p^2+\sigma_{cp}^2+\sigma^2)$ and the variance of
a cell mean over $G$ prompts $\times\,m$ seeds is $\tfrac{\sigma_y^2}{n}[1+(m-1)\rho_{\mathrm{ICC}}]$.
The bracket is the design effect; the effective sample size is
\begin{equation}
n_{\text{eff}}=\frac{n}{1+(m-1)\rho_{\mathrm{ICC}}}.
\label{eq:neff}
\end{equation}
With $m{=}5$ and $\rho_{\mathrm{ICC}}{=}0.3$, $\mathrm{DEFF}{=}2.2$, so 200 generations count as
$\approx91$ independent points---which the MixedLM recovers by partitioning variance. The paired
contrast $d_p=\bar y_{\FULL,p}-\bar y_{\LOO,p}$ cancels $u_p$, which is why cells are seed-matched
within a prompt and the bootstrap resamples whole prompts.

\paragraph{C.4 Kappa paradox, in exact numbers (T4).} General chance-corrected form
$(p_o-p_e)/(1-p_e)$; the coefficients differ only in $p_e$ (Eq.~\ref{eq:agree}). Two raters, 100 binary
items, agreement held at $p_o{=}0.90$. \emph{Skewed} (both-1 $=85$, both-0 $=5$, disagree $=10$):
$\pi_1{=}0.90$, so $p_e{=}0.82\Rightarrow\alphaK{=}1-0.10/0.18{=}0.444$, while
$p_e^{\gamma}{=}0.18\Rightarrow\ac{=}(0.90-0.18)/0.82{=}0.878$. \emph{Balanced} ($\pi_1{=}0.5$):
$p_e{=}p_e^{\gamma}{=}0.5\Rightarrow\alphaK{=}\ac{=}0.80$. Identical agreement, $\alphaK$ moved only by
the marginal---the justification for the disjunctive gate.

\paragraph{C.5 Style-controlled BT and omitted-variable bias (T5).} With
$P(i\succ j)=\sigma((\beta_i-\beta_j)+\gamma^{\top}(s_i-s_j))$, omitting $s$ makes the estimated gap
absorb the cell-mean style shift, $\hat b_i-\hat b_j\to(\beta_i-\beta_j)+\gamma^{\top}(\bar s_i-\bar s_j)$;
if the skill produces longer/denser pages, plain BT inflates the design effect by exactly that term.
Including $s$ identifies $\gamma$ from within-cell style variation and subtracts the bias. The
mediation caveat: colorfulness is partly causal, so it is reported controlled and uncontrolled and its
legitimate effect recovered through the objective channel.

\paragraph{C.6 Paired power (T6).} Eq.~\ref{eq:power} reproduces $\approx$194, 85, 782 decisive pairs
for $60/40$, $65/35$, $55/45$ (rounded to 200/80/780). Total pairs relate to decisive pairs by
$N=N_{\mathrm{dec}}/p_d$; we do not plug the marginal $\delta$ with $p_d{<}1$ into the two-parameter
formula (that mixes a conditional split with a discordance rate). The two-signal rule as error algebra:
the necessity/sufficiency \emph{screen} is $A\cup B$ with $\Pr(\text{false}\mid H_0)\le
\alpha_{\text{obj}}+\alpha_{\text{pref}}\lesssim0.10$; the Stage-2 headline is $E_1\cap E_2\cap E_3\cap
E_4$ with joint false-positive rate $\le\min_k\Pr(E_k\mid H_0)$, far below any single $\alpha$.

\paragraph{C.7 Objective metrics (T7).} WCAG contrast is Eq.~\ref{eq:wcag}. Ngo balance
$\mathrm{BM}=1-\tfrac12(|BM_v|+|BM_h|)$ from summed weight$\times$distance on each side of the frame
axis; equilibrium $\mathrm{EM}=1-\tfrac12(|EM_x|+|EM_y|)$ from the center-of-mass offset; alignment
regularity $=1-n_{\text{distinct edges}}/n_{\text{elements}}$. Hasler colorfulness
$M=\sigma_{rgyb}+0.3\,\mu_{rgyb}$ on opponent axes $rg{=}R-G$, $yb{=}\tfrac12(R+G)-B$. Modular type-scale
adherence is the fraction of adjacent size ratios within $\pm10\%$ (log space) of the best-fitting
canonical ratio. PSI is Eq.~\ref{eq:psi}; POC is Eq.~\ref{eq:poc}.

\paragraph{C.8 Diff-in-means vs.\ LDA, and the ablation projector (T8).} For equal-covariance Gaussians
the Bayes-optimal discriminant is $w^{*}=\Sigma^{-1}(\mu_1-\mu_0)$; diff-in-means $v=\mu_1-\mu_0$ equals
$w^{*}$ iff $\Sigma\propto I$. The whitened direction is optimal for reading but amplifies low-variance
nuisance directions when \emph{writing}, pushing activations off-distribution; diff-in-means is the
on-manifold mean shift and cancels shared nuisance content, so we locate with probes/LDA and steer with
diff-in-means. The ablation operator $P=I-\vhat\vhat^{\top}$ is an orthogonal projector: idempotent
($P^2=I-2\vhat\vhat^{\top}+\vhat(\vhat^{\top}\vhat)\vhat^{\top}=P$ using $\vhat^{\top}\vhat=1$) and
symmetric ($P^{\top}=P$); replacing every residual-writing matrix by $PW_{\text{out}}$ bakes the
ablation into the weights, equivalent to a runtime hook for all writes after the edit.

\paragraph{C.9 Interventions (T9).} Steering and ablation are $\doop$-operations
(Eqs.~\ref{eq:suff}--\ref{eq:flip}); sufficiency uses the \NOSYS{} base so the effect is the direction's,
not a prompt's. The specificity requirement
$\mathbb{E}[Q\mid\doop(h_{\lstar}\mathrel{+}{=}\alpha^{*}\hat r),\NOSYS]\not>\mathbb{E}[Q\mid\NOSYS]$ for
random $\hat r$ of equal norm ($k{=}5$) separates a direction effect from a norm effect. A causal effect
that holds on \emph{unseen task types} adds transportability; the anti-steerable fraction and
dose-response non-monotonicity are reported because a positive mean can hide a heavy tail.

% ------------------------------------------------------------------
\section{Judge template (verbatim)}
\label{app:judge}

\paragraph{Role.} ``You are a senior product designer evaluating two rendered web UIs built for the
\emph{same brief}. Judge overall design quality for that brief. Ignore which looks longer or more
colorful per se; judge whether choices are deliberate and fit the brief.''

\paragraph{Inputs.} the brief text; screenshot A; screenshot B (both desktop $1440\times900$).

\paragraph{Checklist anchors} (from the six metric families, to force reliability): color/palette
coherence; layout \& visual hierarchy; typography; component/pattern quality; restraint (absence of
generic ``AI-slop'' tells); overall fit-to-brief.

\paragraph{Output schema.}
\texttt{\{"winner":"A"|"B"|"tie", "confidence":1-5, "per\_criterion":\{"color":..., "layout":...,
"typography":..., "components":..., "restraint":..., "fit":...\}, "rationale":"$\le$40 words"\}}.

\paragraph{Protocol.} Each pair is judged as (A,B) and (B,A); a winner counts only if consistent across
orders, else tie (Davidson tie term in BT). Judge temperature 0; rubric cached as a system prompt;
batched with backoff. Ties feed the tie term, not a half-win/half-loss split.

% ------------------------------------------------------------------
\section{Additional figures and tables}
\label{app:extra}

Figure~\ref{fig:F4} is the metric-redundancy diagnostic; the POC $z$-scores and averages a preregistered
subset, not all metrics, so redundancy within a family does not double-count. The localization (F5),
stability (F13), PSI-validation (F10), and reliability (F11) figures appear in the main text
(Sections~\ref{sec:oracle}, \ref{sec:stage2method}). Table~T10 (planned vs.\ achieved decisive pairs and
BT-utility simulation power per contrast) is generated from the power module on the observed tie rates.

\synthfig[0.53\linewidth]{F4}{F4_preview.pdf}{4}{Metric-correlation matrix (Spearman $\rho$), blocked by
family. Exposes redundancy so the POC's preregistered subset is defensible; the family-block structure
(within-family redundancy, cross-family separation) is what the real 4{,}630-render matrix shows.
Planted block structure here.}

% ------------------------------------------------------------------
\section{Preregistration summary and amendment log}
\label{app:prereg}

The confirmatory design is frozen before any GPU run and git-tagged \texttt{prereg-v1}. Frozen items:
the 24-cell table (Table~\ref{tab:cells}); the 58-brief split by id (Appendix~\ref{app:corpus}); the POC
formula (Eq.~\ref{eq:poc}) and its two conditionals (PSI-out 6-term variant; $c^{*}$ fallback); the
style-controlled BT endpoint (Eq.~\ref{eq:bt}); the 12-test confirmatory family
$\{$\FULL$-$\NEUTRAL; $5\times$\FULL$-$\LOO-$C_i$; $5\times$\AOI-$C_i-$\NEUTRAL; \FULL$-$\BEAUTY$\}$,
Holm at FWER $0.05$ separately on POC and BT, with everything else BH at FDR $0.10$; the reliability gate
($\alphaK\ge0.667$ or $\ac\ge0.80$ under demonstrable skew); the power gate (achieved $\ge$ required
decisive pairs at the observed split); the Stage-2 success criterion (Section~\ref{sec:stage2method});
and the null rule (Section~\ref{sec:oracle}). The cell table, confirmatory family, success criterion, and
null rule may not be weakened post-freeze; a confirmatory change forks \texttt{prereg-v2} and is reported
as a deviation.

\paragraph{Amendment log.} \emph{(none --- freeze pending the \texttt{prereg-v1} tag).} Post-freeze
amendments will be appended here, dated and numbered, each stating what changed, why, and whether it
touches a confirmatory decision.

% ------------------------------------------------------------------
\section{Data-science lifecycle map}
\label{app:lifecycle}

The paper doubles as the report on each lifecycle phase; Table~\ref{tab:lifecycle} maps section to phase,
to the notebook that carries the narrative, and to the runbook session that produces the real numbers.

\begin{table}[htbp]
\centering\footnotesize
\caption{Paper section $\leftrightarrow$ lifecycle phase $\leftrightarrow$ notebook $\leftrightarrow$
runbook session.}
\label{tab:lifecycle}
\begin{tabular}{@{}p{0.26\linewidth}p{0.20\linewidth}p{0.28\linewidth}p{0.18\linewidth}@{}}
\toprule
paper section & lifecycle phase & notebook & runbook session \\
\midrule
\S\ref{sec:design} study design & problem framing + data & 00--01 & --- (frozen pre-GPU) \\
\S\ref{sec:oracle} oracle & measurement & 02--03 & 4--5 (metrics, judging) \\
\S\ref{sec:stage1} Stage-1 results & experiments + evaluation & 04--05 & 1--3 (pilot, generation) \\
\S\ref{sec:stage2method} Stage-2 method & modeling & 06a/06b, 07a & 6--7 (extract, sweep, freeze) \\
\S\ref{sec:stage2results} Stage-2 results & evaluation & 07a/07b & 8--9 (verify, stretch) \\
\S\ref{sec:discussion}--\ref{sec:repro} & communication + reproducibility & 08 & --- (write-up) \\
\bottomrule
\end{tabular}
\end{table}


\end{document}
